% archon:covers AlgebraicJacobian/Albanese/CodimOneExtension.lean

\chapter{Codimension-1 indeterminacy extension (A.4.a)}
\label{chap:Albanese_CodimOneExtension}

This chapter isolates two codimension-one inputs to the extension theorem for rational
maps.  A rational map from a normal variety to a complete variety is defined away from a
closed subset of codimension at least two, while a map to a group variety has indeterminacy
that is either empty or purely codimension one.  We also record the equivalent
Weil-divisor obstruction in terms of orders along prime divisors.
The depth-\(\geq 2\) reduction uses
\cref{cor:regular_cohen_macaulay} from
\cref{chap:Albanese_AuslanderBuchsbaum}, and the two codimension-one outputs
feed the extension theorem in \cref{chap:Albanese_Thm32RationalMapExtension}.

\section{Setup and motivation}
\label{sec:codim1_setup}

Fix an algebraically closed field \(\bar k\) throughout. By a \emph{variety} over \(\bar k\)
we mean a geometrically integral, separated scheme of finite type over \(\bar k\)
(Milne's convention, \emph{Abelian Varieties}, p.~v); a variety is \emph{nonsingular}
(or \emph{smooth}) if its local ring at every closed point is regular. A variety is
\emph{complete} if it is proper over \(\bar k\).

A \emph{rational map} \(f \colon X \dashrightarrow Y\) between varieties is the equivalence
class of a regular morphism \(\varphi_U \colon U \to Y\) defined on a dense open subset
\(U \subseteq X\), where two pairs \((U, \varphi_U)\), \((U', \varphi_{U'})\) are equivalent
if \(\varphi_U\) and \(\varphi_{U'}\) agree on \(U \cap U'\) (Milne~\S I.3, p.~16) \cite{milne-av}. The map \(f\)
is \emph{defined at} a point \(x \in X\) if there is some representative \((U, \varphi_U)\)
with \(x \in U\); the set of such \(x\) is an open subset \(\mathrm{Dom}(f) \subseteq X\),
the \emph{domain of definition} of \(f\).

The motivating question of this chapter, following Milne~\S I.3, is:
\begin{quote}
  \emph{When does a rational map \(f \colon X \dashrightarrow Y\) extend to a regular
  morphism on all of \(X\)?}
\end{quote}
Without hypotheses on \(X\) or \(Y\) the question has no positive answer (Milne lists three
counterexamples on p.~16: the inclusion of a proper open into the ambient variety; the
rational map to the cuspidal cubic whose inverse fails to extend; the inverse of a
blow-up at a smooth point). The two structural inputs that buy unconditional extension
in our setting are:
\begin{itemize}
  \item \emph{Target completeness:} if \(Y\) is complete (proper over \(\bar k\)), the
    valuative criterion of properness forces the rational map to extend across every
    codim-\(1\) point of a normal source \(X\) (Milne Theorem~3.1).
  \item \emph{Target group structure:} if \(Y\) is a group variety, the difference map
    \(\Phi(x, y) := f(x) \cdot f(y)^{-1}\) on \(X \times X\) has a vanishing/pole divisor
    along the diagonal that controls the codimension of the indeterminacy locus
    (Milne Lemma~3.3).
\end{itemize}
For \(Y\) an abelian variety both conditions hold, and combining them yields Milne's
Theorem~3.2 (the headline of \cref{chap:Albanese_Thm32RationalMapExtension}). The
present chapter packages the two inputs separately:
\cref{thm:codim_one_extension} is Milne~3.1, and
\cref{lem:milne_codim1_indeterminacy} is Milne~3.3. The Weil-divisor reformulation
\cref{thm:weil_divisor_obstruction} translates the codim-\(1\) extension obstruction
in Milne~3.1's proof into the order-at-a-prime-divisor language pinned in
\cref{chap:RiemannRoch_WeilDivisor}.

\section{The indeterminacy locus of a rational map}
\label{sec:indeterm_locus}

\begin{definition}[Indeterminacy locus of a rational map]\leanok
  \label{def:indeterminacy_locus}
  \dcref{custom}
  \lean{AlgebraicGeometry.Scheme.RationalMap.indeterminacyLocus}
  % SOURCE: \cite{milne-av}, \S I.3, p.~16, the paragraph "Let \(V\) and \(W\)
  % be varieties over \(k\)\ldots" defining the domain of definition of a rational map
  % (read from references/abelian-varieties.pdf, PDF page 22).
  % SOURCE QUOTE: "A rational map \(\varphi\) is said to be \emph{defined} at a point
  % \(v\) of \(V\) if \(v \in U\) for some \((U, \varphi_U) \in \varphi\). The set \(U_1\) of
  % \(v\) at which \(\varphi\) is defined is open, and there is a regular map
  % \(\varphi_1 \colon U_1 \to W\) such that \((U_1, \varphi_1) \in \varphi\) --- clearly,
  % \(U_1 = \bigcup_{(U, \varphi_U) \in \varphi} U\) and we can define \(\varphi_1\) to be
  % the regular map such that \(\varphi_1 | U = \varphi_U\) for all $(U, \varphi_U) \in
  % \varphi$."
  \textit{Source: Milne, \emph{Abelian Varieties}, \S I.3, p.~16.}
  Let \(f \colon X \dashrightarrow Y\) be a rational map between varieties over \(\bar k\).
  The \emph{domain of definition} of \(f\) is the open set
  \[
    \mathrm{Dom}(f) \;:=\; \bigcup_{(U, \varphi_U) \in f} U \;\subseteq\; X,
  \]
  i.e.\ the union of the dense opens of all representatives of \(f\). The
  \emph{indeterminacy locus} of \(f\) is the closed complement
  \[
    Z(f) \;:=\; X \smallsetminus \mathrm{Dom}(f) \;\subseteq\; X.
  \]
  The map \(f\) is regular (i.e.\ defined on all of \(X\)) if and only if \(Z(f) = \emptyset\).
\end{definition}

\begin{lemma}
[Indeterminacy locus is closed]
  \label{lem:isClosed_indeterminacy_locus}
  \dcref{custom}
  \uses{def:indeterminacy_locus}
  For any rational map \(f \colon X \dashrightarrow Y\) of schemes, the
  indeterminacy locus \(Z(f) = X \smallsetminus \mathrm{Dom}(f)\)
  (\cref{def:indeterminacy_locus}) is a closed subset of \(X\): it is the
  complement of the open domain of definition \(\mathrm{Dom}(f)\).
\end{lemma}

\begin{proof}
  The domain of definition is open, and the complement of an open subset is closed.
\end{proof}

\begin{definition}[Codimension-1-indeterminacy-free rational maps]\leanok
  \label{def:codim_one_indeterminacy}
  \dcref{custom}
  \lean{AlgebraicGeometry.Scheme.RationalMap.CodimOneFree}
  \uses{def:indeterminacy_locus, def:prime_divisor}
  % SOURCE: \cite{milne-av}, Theorem~3.1, \S I.3, p.~16 --- the conclusion
  % "an open subset \(U\) of \(V\) whose complement \(V \smallsetminus U\) has codimension \(\geq 2\)"
  % is the standing predicate we abstract (read from references/abelian-varieties.pdf,
  % PDF page 22).
  % SOURCE QUOTE: "A rational map \(\varphi \colon V \dashrightarrow W\) from a normal
  % variety \(V\) to a complete variety \(W\) is defined on an open subset \(U\) of \(V\) whose
  % complement \(V \smallsetminus U\) has codimension \(\geq 2\)."
  \textit{Source: Milne, \emph{Abelian Varieties}, Theorem~3.1, \S I.3, p.~16 (the
  conclusion).}
  Let \(f \colon X \dashrightarrow Y\) be a rational map between varieties over \(\bar k\).
  We say \(f\) is \emph{codim-1-indeterminacy-free} (or has \emph{codimension-\(\geq 2\)
  indeterminacy}) if every irreducible component of the closed subset \(Z(f) \subseteq X\)
  has codimension \(\geq 2\) in \(X\). Equivalently, no prime divisor of \(X\) is contained in
  \(Z(f)\); equivalently, every point \(\eta \in X\) with \(\mathrm{codim}_X \overline{\{\eta\}} = 1\)
  lies in \(\mathrm{Dom}(f)\).

  \paragraph{Why the abstraction.} Milne's Theorem~3.1 (\cref{thm:codim_one_extension}
  below) and Lemma~3.3 (\cref{lem:milne_codim1_indeterminacy} below) both produce
  conclusions about the codimension of \(Z(f)\). Phrasing the codim-\(\geq 2\) condition
  as a named predicate keeps the cited statements concise and lets the consumer
  \cref{chap:Albanese_Thm32RationalMapExtension} read off the combination
  (\(Z(f)\) has codim \(\geq 2\) from \cref{thm:codim_one_extension} \emph{and} \(Z(f)\) is
  empty or pure codim 1 from \cref{lem:milne_codim1_indeterminacy} \(\Rightarrow\) \(Z(f)\)
  is empty) as a single line.
\end{definition}

\section{Smoothness yields a DVR at every codim-1 point}
\label{sec:smooth_dvr}

The codim-\(1\) extension argument in Milne's proof of Theorem~3.1 hinges on the local
ring at the generic point of a prime divisor being a discrete valuation ring (so that
the valuative criterion of properness applies). On a normal variety this is the
defining property of regularity in codimension one; on a smooth variety it is automatic.

The smooth-implies-regular-local-stalk implication itself is the Jacobian-criterion
direction of Stacks tag \texttt{00TT}; we package it as a separate sub-lemma so that
\cref{lem:smooth_codim_one_dvr} below reads cleanly as ``regular local of dimension
\(1\) is a DVR'', with the regularity prerequisite supplied by the sub-lemma.

\subsection{K\"ahler differentials under localisation}
\label{subsec:kaehler_localisation_substrate}

The Jacobian criterion for smooth algebras is local.  Near a point of a smooth
scheme one chooses a standard-smooth affine presentation, identifies the stalk
with a localisation of its section ring, and localises the module of K\"ahler
differentials.  The next lemmas record these reductions and the preservation of
freeness and rank under localisation.  The dimension and cotangent-space
comparisons needed for regularity are treated in
\cref{subsec:stage6_subgap_decomposition}.

\begin{lemma}[Smooth implies flatness on stalks]\leanok
  \label{lem:stalkMap_flat_of_smooth}
  \dcref{custom}
  \lean{AlgebraicGeometry.Scheme.stalkMap_flat_of_smooth}
  Let \(X \to \Spec(\bar k)\) be a smooth morphism over an algebraically closed
  field \(\bar k\). Then for every point \(z \in X\) the induced ring map on
  stalks is flat.
\end{lemma}

\begin{proof}
  Every smooth morphism is flat, and flatness of a morphism is equivalent to
  flatness of all induced maps on local rings.
\end{proof}

\begin{lemma}[A smooth morphism is locally standard smooth]\leanok
  \label{lem:exists_standardSmooth_at_of_smooth}
  \dcref{custom}
  \lean{AlgebraicGeometry.Scheme.exists_isStandardSmooth_at_of_smooth}
  For a smooth morphism \(X \to \Spec(\bar k)\) and a point \(z \in X\), there
  exist affine open neighbourhoods \(U \subseteq \Spec(\bar k)\) and \(V \ni z\)
  in \(X\) with \(V \subseteq X^{-1}(U)\) such that the induced ring map
  \(\Gamma(\Spec \bar k, U) \to \Gamma(X, V)\) is standard smooth (Stacks 00T7).
\end{lemma}

\begin{proof}
  This is the local standard-smooth presentation theorem for smooth morphisms
  (Stacks tag~00T7).
\end{proof}

\begin{lemma}[A standard-smooth presentation of a smooth stalk]\leanok
  \label{lem:exists_algebra_standardSmooth_stalk_localization}
  \dcref{custom}
  \lean{AlgebraicGeometry.Scheme.exists_algebra_isStandardSmooth_section_stalk_isLocalization_of_smooth}
  \uses{lem:exists_standardSmooth_at_of_smooth, lem:isLocalization_atPrime_stalk_of_affineOpen}
  For a smooth morphism \(X \to \Spec(\bar k)\) and a point \(z \in X\), there is
  an affine open neighbourhood \(V \ni z\) and an \(\Gamma(\Spec \bar k, U)\)-algebra
  structure on \(\Gamma(X, V)\) making it a standard-smooth algebra, together with a
  witness that the stalk \(\mathcal O_{X, z}\) is the localisation of \(\Gamma(X, V)\)
  at the prime ideal corresponding to \(z\).
\end{lemma}

\begin{proof}
  Choose the affine neighbourhoods supplied by
  \cref{lem:exists_standardSmooth_at_of_smooth}.  Their map on section rings is
  standard smooth, and
  \cref{lem:isLocalization_atPrime_stalk_of_affineOpen} identifies the stalk at
  \(z\) with the localisation of the target section ring at the prime of \(z\).
\end{proof}

\begin{lemma}[K\"ahler differentials of a standard-smooth algebra are free]\leanok
  \label{lem:module_free_kaehler_standardSmooth}
  \dcref{custom}
  \lean{AlgebraicGeometry.Scheme.module_free_kaehlerDifferential_of_isStandardSmooth}
  For a standard-smooth \(R\)-algebra \(S\), the module of K\"ahler differentials
  \(\Omega_{S/R}\) is free as an \(S\)-module (Stacks 00T7(2)), with basis the family
  \(\{ds_i\}\) of the standard-smooth presentation.
\end{lemma}

\begin{proof}
  The invertible Jacobian minor of a standard-smooth presentation makes the
  differentials of the free variables a basis of \(\Omega_{S/R}\).
\end{proof}

\begin{lemma}[K\"ahler-differential freeness is preserved by localisation]\leanok
  \label{lem:module_free_kaehler_localization}
  \dcref{custom}
  \lean{AlgebraicGeometry.Scheme.module_free_kaehlerDifferential_localization}
  \uses{lem:module_free_kaehler_standardSmooth}
  Let \(R\) and \(S\) be commutative rings equipped with an \(R\)-algebra
  structure on \(S\), and assume that the module of K\"ahler differentials
  \(\Omega_{S/R}\) is free as an \(S\)-module. Let \(M \subseteq S\) be a
  submonoid and let \(S_M\) be a localisation of \(S\) at \(M\), so that
  \(R \to S \to S_M\) is a tower of \(R\)-algebras. Then the module of
  K\"ahler differentials \(\Omega_{S_M / R}\) is free as an
  \(S_M\)-module.
\end{lemma}

\begin{proof}
  The canonical comparison map
  \(S_M \otimes_S \Omega_{S/R} \to \Omega_{S_M/R}\) is an isomorphism.
  Localising a basis of \(\Omega_{S/R}\) therefore gives a basis of
  \(\Omega_{S_M/R}\) over \(S_M\).
\end{proof}

\begin{lemma}[The rank of localised K\"ahler differentials equals the relative dimension]\leanok
  \label{lem:rank_kaehler_localization_eq_relative_dim}
  \dcref{custom}
  \lean{AlgebraicGeometry.Scheme.rank_kaehlerDifferential_localization_eq_relativeDimension}
  \uses{lem:module_free_kaehler_localization}
  With notation as in \cref{lem:module_free_kaehler_localization},
  assume in addition that the \(R\)-algebra \(S\) is \emph{standard
  smooth of relative dimension \(n\)} (in the sense of Stacks tag
  \texttt{00T7}(2)) and that
  the localisation \(S_M\) is non-trivial as a ring. Then the rank of
  \(\Omega_{S_M / R}\) over \(S_M\) equals \(n\).
\end{lemma}

\begin{proof}
  A standard-smooth presentation of relative dimension \(n\) gives
  \(\Omega_{S/R}\) a basis indexed by an \(n\)-element set.  By
  \cref{lem:module_free_kaehler_localization}, localising this basis gives
  a basis of \(\Omega_{S_M/R}\) with the same index set.  Its rank is
  therefore \(n\).
\end{proof}

\subsection{Dimension and cotangent criteria for smooth local rings}
\label{subsec:stage6_subgap_decomposition}

Let \(k\) be a perfect field and let \(k \to S\) be a standard-smooth
presentation near a point \(z\),
and write \(S_{\mathfrak q}=\mathcal O_{X,z}\) for the corresponding
localisation.  The regularity of this local ring follows by comparing
three dimensions.  The finite-type dimension formula (Stacks tag~00OE)
and the cotangent exact sequence (Stacks tag~02JK) give
\[
  \dim S_{\mathfrak q}
  = n-\operatorname{trdeg}_k\kappa(\mathfrak q)
  = \dim_{\kappa(\mathfrak q)} \mathfrak m / \mathfrak m^2,
\]
where \(n=\operatorname{rank}_{S_{\mathfrak q}}
\Omega_{S_{\mathfrak q}/k}\),
which is the cotangent-space characterisation of a regular local ring.
The next three lemmas state the dimension formula, the cotangent
comparison, and their combination.

\begin{lemma}
  [Smooth-algebra Krull-dimension formula (Stacks 00OE)]
  \label{lem:smooth_algebra_krull_dim_formula}
  \dcref{custom}
  \lean{Algebra.IsStandardSmoothOfRelativeDimension.ringKrullDim_localization_eq_relativeDimension}
  \uses{lem:rank_kaehler_localization_eq_relative_dim,
        lem:ringKrullDim_localization_atMaximal_mvPolynomial,
        lem:ringKrullDim_quotient_localization_mvPolynomial_regular,
        lem:matsumura_isRegular_of_linearIndependent_cotangent,
        lem:exists_submersivePresentation_of_standardSmoothReldim,
        lem:submersive_relation_cotangent_linearIndependent_localized}
  % SOURCE: \cite{stacks-project}, lemma-dimension-at-a-point-finite-type-field
  % (canonically known as tag 00OE) (read from
  % references/stacks-algebra.tex, L28206--28217 in section
  % "Dimension theory of finitely generated $k$-algebras").
  % SOURCE QUOTE: "Let $k$ be a field. Let $S$ be a finite type $k$
  % algebra. Let $X = \Spec(S)$. Let $\mathfrak p \subset S$ be a prime
  % ideal, and let $x \in X$ be the corresponding point. Then we have
  % $\dim_x(X) = \dim(S_{\mathfrak p}) + \text{trdeg}_k\
  % \kappa(\mathfrak p)$."
  \textit{Source: \cite{stacks-project}, tag~00OE (\(=\) \texttt{lemma-dimension-at-a-point-finite-type-field}).}
  Let \(k\) be a field, let \(S\) be a finite-type \(k\)-algebra, and let
  \(\mathfrak q \subset S\) be a prime ideal corresponding to a point
  \(x \in \mathrm{Spec}(S)\). Then
  \[
    \dim_x(\mathrm{Spec}\,S)
    \;=\; \dim S_{\mathfrak q} + \mathrm{trdeg}_k\,\kappa(\mathfrak q).
  \]
  In particular, if \(\mathfrak q\) is a maximal ideal (so \(x\) is a
  closed point) and \(k\) is algebraically closed, then
  \(\kappa(\mathfrak q) = k\) (Nullstellensatz), so
  \(\mathrm{trdeg}_k\,\kappa(\mathfrak q) = 0\) and
  \(\dim S_{\mathfrak q} = \dim_x(\mathrm{Spec}\,S)\).

  \paragraph{Geometric application.} Let \(k = \bar k\), let
  \(S\) be a standard-smooth \(\bar k\)-algebra of relative dimension
  \(n\), and let \(\mathfrak q\) be the prime defining \(z \in \Spec S\).  The
  formula yields \(\dim(\mathcal O_{X, z}) = \dim S_{\mathfrak q} = n -
  \mathrm{trdeg}_{\bar k}\,\kappa(z)\). For \(z\) a closed point this is
  simply \(n\); for the generic point of a codim-\(1\) subvariety
  (\(\mathrm{trdeg} = n - 1\)) it is \(1\).
\end{lemma}

% SOURCE QUOTE PROOF: "By Lemma \ref{lemma-dimension-prime-polynomial-ring}
% we know that $r = \text{trdeg}_k\ \kappa(\mathfrak p)$ is equal to the
% dimension of $V(\mathfrak p)$. Pick any maximal chain of primes
% $\mathfrak p \subset \mathfrak p_1 \subset \ldots \subset \mathfrak p_r$
% starting with $\mathfrak p$ in $S$. This has length $r$ by Lemma
% \ref{lemma-dimension-spell-it-out}. Let $\mathfrak q_j$, $j \in J$ be
% the minimal primes of $S$ which are contained in $\mathfrak p$. These
% correspond $1-1$ to minimal primes in $S_{\mathfrak p}$ via the rule
% $\mathfrak q_j \mapsto \mathfrak q_jS_{\mathfrak p}$. By Lemma
% \ref{lemma-dimension-at-a-point-finite-type-over-field} we know that
% $\dim_x(X)$ is equal to the maximum of the dimensions of the rings
% $S/\mathfrak q_j$. For each $j$ pick a maximal chain of primes
% $\mathfrak q_j \subset \mathfrak p'_1 \subset \ldots \subset
% \mathfrak p'_{s(j)} = \mathfrak p$. Then $\dim(S_{\mathfrak p}) =
% \max_{j \in J} s(j)$. Now, each chain
% $\mathfrak q_j \subset \mathfrak p'_1 \subset \ldots \subset
% \mathfrak p'_{s(j)} = \mathfrak p \subset \mathfrak p_1 \subset \ldots
% \subset \mathfrak p_r$ is a maximal chain in $S/\mathfrak q_j$, and by
% what was said before we have $\dim_x(X) = \max_{j \in J} r + s(j)$.
% The lemma follows."
\begin{proof}
  \uses{lem:rank_kaehler_localization_eq_relative_dim,
        lem:ringKrullDim_localization_atMaximal_mvPolynomial,
        lem:ringKrullDim_quotient_localization_mvPolynomial_regular,
        lem:matsumura_isRegular_of_linearIndependent_cotangent,
        lem:exists_submersivePresentation_of_standardSmoothReldim,
        lem:submersive_relation_cotangent_linearIndependent_localized}
  This is exactly Stacks tag~\texttt{00OE}
  (\texttt{lemma-dimension-at-a-point-finite-type-field}); the proof is a
  citation of the dimension theory of finitely generated algebras over a
  field. By Noether normalisation the finite-type \(k\)-algebra \(S\) is
  finite over a polynomial subalgebra, so the going-up / catenary
  dimension theory of such algebras applies. Write
  \(r = \mathrm{trdeg}_k\,\kappa(\mathfrak q)\); this transcendence
  degree is exactly the dimension of the closed subvariety
  \(\overline{\{x\}} = V(\mathfrak q)\). A maximal chain of primes
  through \(\mathfrak q\) therefore splits into a
  length-\(\dim S_{\mathfrak q}\) part below \(\mathfrak q\) (recording
  the local codimension, i.e.\ the height of \(\mathfrak q\) in the
  relevant irreducible component) and a length-\(r\) part above it
  (recording the dimension of \(\overline{\{x\}}\)). Taking the maximum
  over the minimal primes contained in \(\mathfrak q\) yields
  \[
    \dim_x(\mathrm{Spec}\,S) = \dim S_{\mathfrak q} + r
    = \dim S_{\mathfrak q} + \mathrm{trdeg}_k\,\kappa(\mathfrak q),
  \]
  which is the asserted formula.

  For the geometric specialisation, take
  \(k = \bar k\) algebraically closed and \(S\) standard smooth of
  relative dimension \(n\), so that \(\dim_x(\mathrm{Spec}\,S) = n\) via
  \cref{lem:rank_kaehler_localization_eq_relative_dim} (the free
  K\"ahler module \(\Omega_{S_{\mathfrak q}/R}\) has rank \(n\), which
  pins the relative dimension). The formula then gives
  \(\dim S_{\mathfrak q} = n - \mathrm{trdeg}_{\bar k}\,\kappa(\mathfrak
  q)\). At a closed point \(\mathfrak q\) is maximal, so by the
  Nullstellensatz \(\kappa(\mathfrak q) = \bar k\) and
  \(\mathrm{trdeg}_{\bar k}\,\kappa(\mathfrak q) = 0\); hence
  \(\dim S_{\mathfrak q} = \dim_x(\mathrm{Spec}\,S) = n\).
\end{proof}

\begin{lemma}
  [Cotangent and K\"ahler differentials over a field]
  \label{lem:cotangent_kahler_over_field}
  \dcref{custom}
  \lean{AlgebraicGeometry.Scheme.cotangent_iso_residue_tensor_kaehler_of_formallySmooth_residue}
  \uses{lem:smooth_algebra_krull_dim_formula}
  % SOURCE: \cite{stacks-project}, lemma-differential-seq (canonically
  % known as the right-exact conormal portion of tag 02JK; the
  % short-exact / split refinement is the proof step inside
  % lemma-separable-smooth) (read from
  % references/stacks-algebra.tex, L34087--34104 and L38724--38766).
  % SOURCE QUOTE: "In diagram (\ref{equation-functorial-omega}),
  % suppose that $S \to S'$ is surjective with kernel $I \subset S$, and
  % assume that $R' = R$. Then there is a canonical exact sequence of
  % $S'$-modules
  % $$I/I^2 \longrightarrow \Omega_{S/R} \otimes_S S' \longrightarrow
  % \Omega_{S'/R} \longrightarrow 0$$
  % The leftmost map is characterized by the rule that $f \in I$ maps
  % to $\text{d}f \otimes 1$."
  % SOURCE QUOTE: "Hence we get
  % $$\dim_{\kappa} \Omega_{R/k} \otimes_R \kappa = \dim_\kappa
  % \mathfrak m/\mathfrak m^2 + \text{trdeg}_k (\kappa) \geq \dim R +
  % \text{trdeg}_k (\kappa) = \dim_{\mathfrak q} S$$
  % (see Lemma \ref{lemma-dimension-at-a-point-finite-type-field} for
  % the last equality) with equality if and only if $R$ is regular."
  \textit{Source: \cite{stacks-project}, tag~02JK. The right-exact half is
  \texttt{lemma-differential-seq}; the separable-residue refinement occurs in
  \texttt{lemma-separable-smooth}.}
  Let \(k\) be a field and let \(R\) be a Noetherian local
  \(k\)-algebra with maximal ideal \(\mathfrak m\) and residue field
  \(\kappa = R/\mathfrak m\). The canonical exact sequence
  \[
    \mathfrak m / \mathfrak m^2
    \;\longrightarrow\; \Omega_{R/k} \otimes_R \kappa
    \;\longrightarrow\; \Omega_{\kappa / k}
    \;\longrightarrow\; 0
  \]
  identifies the cotangent space \(\mathfrak m / \mathfrak m^2\) (a
  \(\kappa\)-vector space) with the kernel of the right map. When
  \(\kappa / k\) is separable (in particular when \(k\) is
  algebraically closed and \(\kappa = k\), so \(\Omega_{\kappa / k}
  = 0\)) the leftmost map is injective and the sequence is short exact
  --- in fact split short exact when there is a \(k\)-algebra section
  \(\kappa \to R\) (e.g.\ at a closed point of a \(\bar k\)-variety).
  In particular, over an algebraically closed field \(k = \bar k\) and
  at a closed point \(\mathfrak m\),
  \[
    \Omega_{R/k} \otimes_R \kappa \;\simeq\; \mathfrak m / \mathfrak m^2
  \]
  as \(\kappa\)-vector spaces; hence
  \(\dim_\kappa(\Omega_{R/k} \otimes_R \kappa)
   = \dim_\kappa \mathfrak m / \mathfrak m^2\).
\end{lemma}

The cotangent--K\"ahler comparison can be separated into the following
finite-dimensional statements.  They compute the dimension after residue-field
base change and then transfer it across the isomorphism of
\cref{lem:cotangent_kahler_over_field}.

\begin{lemma}[Residue-field tensor dimension from free rank]\leanok
  \label{lem:finrank_residue_kaehler_of_free_rank}
  \dcref{custom}
  \lean{AlgebraicGeometry.Scheme.finrank_residueField_tensor_kaehlerDifferential_of_free_rank_eq}
  Let \(R \to S_{\mathfrak q}\) be an algebra with \(S_{\mathfrak q}\) a nontrivial
  local ring whose K\"ahler differentials \(\Omega_{S_{\mathfrak q}/R}\) are
  \(S_{\mathfrak q}\)-free of rank \(n\). Then the residue-field base change
  \(\kappa \otimes_{S_{\mathfrak q}} \Omega_{S_{\mathfrak q}/R}\) has
  \(\kappa\)-dimension \(n\), where \(\kappa\) is the residue field of \(S_{\mathfrak q}\).
\end{lemma}

\begin{proof}
  A basis of the free module \(\Omega_{S_{\mathfrak q}/R}\) base-changes to a
  basis of \(\kappa \otimes_{S_{\mathfrak q}}\Omega_{S_{\mathfrak q}/R}\)
  with the same \(n\)-element index set.
\end{proof}

\begin{lemma}[Residue-tensor dimension for a standard-smooth algebra]
  \label{lem:finrank_residue_kaehler_of_standardSmooth_reldim}
  \dcref{custom}
  \lean{AlgebraicGeometry.Scheme.finrank_residueField_tensor_kaehlerDifferential_of_isStandardSmoothOfRelativeDimension}
  \uses{lem:finrank_residue_kaehler_of_free_rank}
  If \(S_{\mathfrak q}\) is a nontrivial local \(R\)-algebra that is standard smooth
  of relative dimension \(n\), then the residue-field base change of its K\"ahler
  differentials has \(\kappa\)-dimension \(n\). This is the standard-smooth-tied
  packaging of \cref{lem:finrank_residue_kaehler_of_free_rank}, discharging the
  freeness and rank hypotheses directly from the relative-dimension instance.
\end{lemma}

\begin{proof}
  Standard smoothness of relative dimension \(n\) makes
  \(\Omega_{S_{\mathfrak q}/R}\) free of rank \(n\), so apply
  \cref{lem:finrank_residue_kaehler_of_free_rank}.
\end{proof}

\begin{lemma}[Cotangent isomorphism in maximal-ideal form]
  \label{lem:cotangent_iso_maximalIdeal_residue_kaehler}
  \dcref{custom}
  \lean{AlgebraicGeometry.Scheme.cotangent_iso_maximalIdeal_residue_tensor_kaehler_of_formallySmooth_residue}
  \uses{lem:cotangent_kahler_over_field}
  Under the closed-point hypotheses of \cref{lem:cotangent_kahler_over_field}
  (\(S_{\mathfrak q}\) formally smooth over \(R\), residue field \(\kappa\) formally
  smooth over \(R\), and \(\Omega_{\kappa/R} = 0\)), there is an
  \(S_{\mathfrak q}\)-linear isomorphism between the cotangent module of the maximal
  ideal \(\mathfrak m\) and the residue-field base change
  \(\kappa \otimes_{S_{\mathfrak q}} \Omega_{S_{\mathfrak q}/R}\). This restates the
  kernel-side cotangent isomorphism of \cref{lem:cotangent_kahler_over_field} with the
  canonical maximal-ideal cotangent space as its domain.
\end{lemma}

\begin{proof}
  Transport the isomorphism of
  \cref{lem:cotangent_kahler_over_field} along the canonical equality
  \(\ker(S_{\mathfrak q} \to \kappa) = \mathfrak m\).
\end{proof}

\begin{lemma}[Cotangent-space dimension with formally smooth residue field]\leanok
  \label{lem:finrank_cotangentSpace_of_formallySmooth_residue}
  \dcref{custom}
  \lean{AlgebraicGeometry.Scheme.finrank_cotangentSpace_of_formallySmooth_residue}
  \uses{lem:cotangent_iso_maximalIdeal_residue_kaehler, lem:finrank_residue_kaehler_of_free_rank}
  With the closed-point hypotheses of
  \cref{lem:cotangent_iso_maximalIdeal_residue_kaehler} and
  \(\Omega_{S_{\mathfrak q}/R}\) free of \(S_{\mathfrak q}\)-rank \(n\), the cotangent
  space \(\mathfrak m / \mathfrak m^2\) has \(\kappa\)-dimension \(n\).
\end{lemma}

\begin{proof}
  Combine the cotangent isomorphism
  \cref{lem:cotangent_iso_maximalIdeal_residue_kaehler} with the residue-tensor
  dimension formula \cref{lem:finrank_residue_kaehler_of_free_rank}.
\end{proof}

\begin{lemma}[Cotangent-space dimension at a $\bar k$-rational point]\leanok
  \label{lem:finrank_cotangentSpace_of_bijective_residue}
  \dcref{custom}
  \lean{AlgebraicGeometry.Scheme.finrank_cotangentSpace_of_bijective_algebraMap_residue}
  \uses{lem:finrank_cotangentSpace_of_formallySmooth_residue}
  Bundled closed-point form of
  \cref{lem:finrank_cotangentSpace_of_formallySmooth_residue}: if \(S_{\mathfrak q}\)
  is formally smooth over \(R\), \(\Omega_{S_{\mathfrak q}/R}\) is free of rank \(n\),
  and the structure map \(R \to \kappa\) to the residue field is bijective (as at a
  \(\bar k\)-rational closed point, by the Nullstellensatz), then
  \(\dim_\kappa \mathfrak m/\mathfrak m^2 = n\).
\end{lemma}

\begin{proof}
  A bijective structure map identifies the residue field with \(R\); it is
  therefore formally smooth over \(R\), with zero module of relative
  differentials.  The result follows from
  \cref{lem:finrank_cotangentSpace_of_formallySmooth_residue}.
\end{proof}

\begin{lemma}
  [A standard-smooth stalk is regular local]
  \label{lem:stage6_regular_stalk_assembly}
  \dcref{custom}
  \lean{AlgebraicGeometry.TODO.stage6_regular_stalk_assembly}
  \uses{lem:rank_kaehler_localization_eq_relative_dim,
        lem:smooth_algebra_krull_dim_formula,
        lem:cotangent_kahler_over_field}
  Let \(\bar k\) be an algebraically closed field, and let
  \(\bar k \to S\) be a standard-smooth algebra presentation of
  relative dimension \(n\) on an affine chart \(\Spec S \subseteq X\),
  and let \(z \in \Spec S\) be a point corresponding to a prime
  \(\mathfrak q \subset S\). Then the stalk \(S_{\mathfrak q}\) is a
  regular local ring of Krull dimension
  \(n - \mathrm{trdeg}_{\bar k}\,\kappa(\mathfrak q)\). In the
  closed-point case (\(\kappa(\mathfrak q) = \bar k\),
  \(\mathrm{trdeg} = 0\)) the dimension equals \(n\); in the
  codim-\(1\) generic-point case (\(\mathrm{trdeg} = n - 1\)) the
  dimension equals \(1\).
\end{lemma}

\begin{proof}
  Combine the three sub-lemmas:
  \begin{enumerate}
    \item By \cref{lem:rank_kaehler_localization_eq_relative_dim},
      the localised K\"ahler module \(\Omega_{S_{\mathfrak q} / \bar k}\)
      is free of rank \(n\) over \(S_{\mathfrak q}\).
    \item By \cref{lem:smooth_algebra_krull_dim_formula}, the Krull dimension of
      \(S_{\mathfrak q}\) equals
      \(\dim_z \Spec S - \mathrm{trdeg}_{\bar k}\,\kappa(\mathfrak q)\),
      which for a standard-smooth presentation of relative dimension
      \(n\) equals \(n - \mathrm{trdeg}_{\bar k}\,\kappa(\mathfrak q)\).
      (For the standard-smooth presentation the dimension
      \(\dim_z \Spec S = n\) follows from the global-complete-intersection
      structure: \(n\) variables modulo \(n - \dim\) equations, see
      Stacks \texttt{lemma-relative-global-complete-intersection-smooth}.)
    \item By \cref{lem:cotangent_kahler_over_field}, the cotangent exact
      sequence is short exact because \(\bar k\) is perfect and
      \(\kappa(\mathfrak q)/\bar k\) is separably generated.  Its middle
      term has dimension \(n\) by (1), while
      \(\Omega_{\kappa(\mathfrak q)/\bar k}\) has dimension
      \(\operatorname{trdeg}_{\bar k}\kappa(\mathfrak q)\).  Hence
      \[
        \dim_{\kappa(\mathfrak q)}\mathfrak m/\mathfrak m^2
        =n-\operatorname{trdeg}_{\bar k}\kappa(\mathfrak q).
      \]
      At a closed point the transcendence degree is zero, and the
      cotangent sequence reduces to the displayed isomorphism in
      \cref{lem:cotangent_kahler_over_field}.
    \item Compare (2) and (3): both \(\dim S_{\mathfrak q}\) and
      \(\dim_{\kappa(\mathfrak q)} \mathfrak m / \mathfrak m^2\) equal
      \(n - \mathrm{trdeg}_{\bar k}\,\kappa(\mathfrak q)\); the
      regularity criterion for Noetherian local rings then gives \(S_{\mathfrak q}\)
      regular. The scheme-level transfer
      \(S_{\mathfrak q} = \mathcal O_{X, z}\) on the affine chart completes
      the conclusion.
  \end{enumerate}
\end{proof}

\subsection{Krull dimension of standard-smooth localisations}
\label{subsec:stage6_iib_substrate_iter200}

For a prime \(\mathfrak p\) of a Noetherian ring, the Krull dimension
of \(R_{\mathfrak p}\) is the height of \(\mathfrak p\).  If
\(R=k[x_i]_{i\in\iota}\) is a polynomial ring in finitely many variables
over a field and \(\mathfrak m\) is maximal, then
\(\operatorname{ht}(\mathfrak m)=\#\iota\).  The lower bound follows by
induction on the number of variables and the principal ideal theorem;
the upper bound follows from the dimension of the polynomial ring.
Consequently \(\dim R_{\mathfrak m}=\#\iota\).

If \(f_1,\ldots,f_c\) is a regular sequence in a Noetherian local ring
\(A\), then
\[
  \dim A/(f_1,\ldots,f_c)+c=\dim A.
\]
Thus a standard-smooth presentation gives the desired local dimension
once its defining relations are known to form a regular sequence after
localisation.

For a submersive presentation, the invertible Jacobian minor makes the
cotangent classes of the defining relations linearly independent.  This
independence is preserved after localisation.  Matsumura's criterion then
shows that the localised relations form a regular sequence.  More explicitly:
\begin{itemize}
  \item On a regular local Noetherian ring \((A, \mathfrak m)\), a
    sequence \(f_1, \dots, f_c \in \mathfrak m\) whose images in
    \(\mathfrak m / \mathfrak m^2\) are \(\kappa\)-linearly
    independent forms a regular sequence in \(A\).  Indeed, \(f_1\notin
    \mathfrak m^2\), so it is a non-zero-divisor and \(A/(f_1)\) is again
    regular local; the remaining cotangent classes stay independent in
    the quotient, and induction applies.

  \item The conormal isomorphism
    \[
      (I/I^2) \otimes_S S_{\mathfrak m}
      \simeq (I S_{\mathfrak m})/(I S_{\mathfrak m})^2
    \]
    transports linear independence from \(I/I^2\) over \(S\) to
    \(\mathfrak m_A / \mathfrak m_A^2\) over \(\kappa(\mathfrak m_A)\)
    on the localisation \(A\).
  \item Applying the preceding regular-sequence criterion and the
    dimension-drop formula gives the dimension of the quotient of the
    localised polynomial ring by the relations.
\end{itemize}

\subsection{Algebraic ingredients and scheme bridges}
\label{subsec:codim1_substrate_lemmas}

The following lemmas make the preceding dimension argument explicit. They feed
the Stacks-00OE Krull-dimension formula
\cref{lem:smooth_algebra_krull_dim_formula} and the regular-stalk assembly,
either directly or through the Jacobian regular-sequence witness.

\paragraph{Stacks 00OE Krull-dimension chain.}

\begin{lemma}[Localised Krull dimension equals height]\leanok
  \label{lem:ringKrullDim_localization_eq_height_atPrime}
  \dcref{custom}
  \lean{AlgebraicGeometry.Scheme.ringKrullDim_localization_eq_height_atPrime}
  For a prime ideal \(\mathfrak m\) of a commutative ring \(R\) and any localisation
  \(A\) of \(R\) at \(\mathfrak m\), the Krull dimension of \(A\) equals the height of
  \(\mathfrak m\).
\end{lemma}

\begin{proof}
  Prime ideals of \(R_{\mathfrak m}\) correspond order-preservingly to
  prime ideals of \(R\) contained in \(\mathfrak m\).  Thus chains of
  primes in the localisation have the same maximal lengths as chains
  ending at \(\mathfrak m\), whose supremum is its height.
\end{proof}

\begin{lemma}
[Polynomial-ring maximal ideal height lower bound]
  \label{lem:mvPolynomial_maximalIdeal_height_ge_card}
  \dcref{custom}
  For a field \(k\), every maximal ideal of the polynomial ring
  \(\mathrm{MvPolynomial}\,(\mathrm{Fin}\,n)\,k\) has height at least \(n\).
\end{lemma}

\begin{proof}
  Induct on \(n\), transporting the ideal along the
  one-variable-splitting isomorphism and using that adjoining a polynomial variable
  raises height by one, together with the Jacobson contraction of the maximal ideal.
\end{proof}

\begin{lemma}[Polynomial-ring ideal height upper bound]
  \label{lem:mvPolynomial_maximalIdeal_height_le_natCard}
  \dcref{custom}
  \lean{AlgebraicGeometry.Scheme.MvPolynomial.maximalIdeal_height_le_natCard_of_field}
  For a field \(k\) and a finite index set \(\iota\), every proper ideal of
  \(\mathrm{MvPolynomial}\,\iota\,k\) has height at most \(\#\iota\).
\end{lemma}

\begin{proof}
  The height of a proper ideal is bounded by the Krull dimension of its
  ring, and the polynomial ring in \(\#\iota\) variables over a field has
  Krull dimension \(\#\iota\).
\end{proof}

\begin{lemma}
[Maximal-ideal height in a polynomial ring indexed by $\mathrm{Fin}\,n$]
  \label{lem:mvPolynomial_maximalIdeal_height_eq_card}
  \dcref{custom}
  \uses{lem:mvPolynomial_maximalIdeal_height_ge_card, lem:mvPolynomial_maximalIdeal_height_le_natCard}
  For a field \(k\), every maximal ideal of
  \(\mathrm{MvPolynomial}\,(\mathrm{Fin}\,n)\,k\) has height exactly \(n\).
\end{lemma}

\begin{proof}
  Combine the lower bound
  \cref{lem:mvPolynomial_maximalIdeal_height_ge_card} and the upper bound
  \cref{lem:mvPolynomial_maximalIdeal_height_le_natCard} by antisymmetry.
\end{proof}

\begin{lemma}[Maximal-ideal height in a finitely generated polynomial ring]
  \label{lem:mvPolynomial_maximalIdeal_height_eq_natCard}
  \dcref{custom}
  \lean{AlgebraicGeometry.Scheme.MvPolynomial.maximalIdeal_height_eq_natCard}
  \uses{lem:mvPolynomial_maximalIdeal_height_eq_card}
  For a field \(k\) and a finite index set \(\iota\), every maximal ideal of
  \(\mathrm{MvPolynomial}\,\iota\,k\) has height exactly \(\#\iota\).
\end{lemma}

\begin{proof}
  Choose a bijection \(\iota \simeq \mathrm{Fin}(\#\iota)\) and transport
  \cref{lem:mvPolynomial_maximalIdeal_height_eq_card} along the renaming isomorphism
  that it induces. Ring isomorphisms preserve heights of ideals.
\end{proof}

\begin{lemma}[Krull dimension of a localised polynomial ring]
  \label{lem:ringKrullDim_localization_atMaximal_mvPolynomial}
  \dcref{custom}
  \lean{AlgebraicGeometry.Scheme.ringKrullDim_localization_atMaximal_MvPolynomial}
  \uses{lem:ringKrullDim_localization_eq_height_atPrime, lem:mvPolynomial_maximalIdeal_height_eq_natCard}
  For a field \(k\), a finite index set \(\iota\), and any localisation \(A\) of
  \(\mathrm{MvPolynomial}\,\iota\,k\) at a maximal ideal \(\mathfrak m\), the Krull
  dimension of \(A\) equals \(\#\iota\).
\end{lemma}

\begin{proof}
  By
  \cref{lem:ringKrullDim_localization_eq_height_atPrime} (Krull dimension equals
  height), the dimension is the height of \(\mathfrak m\), which is
  \(\#\iota\) by \cref{lem:mvPolynomial_maximalIdeal_height_eq_natCard}.
\end{proof}

\begin{lemma}
[Dimension drop by a regular sequence]
  \label{lem:ringKrullDim_quotient_add_eq_regular_sequence}
  \dcref{custom}
  Let \(R'\) be a Noetherian local ring of dimension \(a\), and let
  \(rs\) be a regular sequence in \(R'\) of length \(b\). Then
  \(\dim(R'/(rs)) + b = a\).
\end{lemma}

\begin{proof}
  Quotienting a Noetherian local ring by one non-zero-divisor lowers its
  dimension by one.  Apply this successively to the elements of the
  regular sequence.
\end{proof}

\begin{lemma}[Dimension of a localised polynomial quotient]
  \label{lem:ringKrullDim_quotient_localization_mvPolynomial_regular}
  \dcref{custom}
  \lean{AlgebraicGeometry.Scheme.ringKrullDim_quotient_localization_MvPolynomial_of_regular}
  \uses{lem:ringKrullDim_localization_atMaximal_mvPolynomial, lem:ringKrullDim_quotient_add_eq_regular_sequence}
  Let \(A\) be a Noetherian localisation of \(\mathrm{MvPolynomial}\,\iota\,k\) at a
  maximal ideal, and let \(rs\) be a regular sequence in \(A\) of length \(b\). Then
  \(\mathrm{ringKrullDim}\,(A/(rs)) + b = \#\iota\).
\end{lemma}

\begin{proof}
  The localised polynomial ring has dimension \(\#\iota\) by
  \cref{lem:ringKrullDim_localization_atMaximal_mvPolynomial}; now apply
  \cref{lem:ringKrullDim_quotient_add_eq_regular_sequence}.
\end{proof}

\paragraph{Stacks 00SW submersive-presentation cotangent independence.}

\begin{lemma}
[00SW --- relations are cotangent-independent]
  \label{lem:submersive_relation_cotangent_linearIndependent}
  \dcref{custom}
  For a submersive presentation \(P\) of an \(R\)-algebra \(S\), the family of
  cotangent classes of the relations \(i \mapsto [\,P.\mathrm{relation}\,i\,]\) in the
  cotangent module of \(P\) is \(S\)-linearly independent.
\end{lemma}

\begin{proof}
  In a submersive presentation, the cotangent classes of the relations
  are the distinguished cotangent basis, hence are linearly independent.
\end{proof}

\begin{lemma}
[00SW --- cotangent independence transported to a localisation]
  \label{lem:submersive_relation_cotangent_linearIndependent_localized}
  \dcref{custom}
  \uses{lem:submersive_relation_cotangent_linearIndependent}
  With \(P\) a submersive presentation of \(S\) and \(p\) a prime ideal of \(S\), the
  image of the cotangent classes of the relations under the canonical localisation map
  of the cotangent module at \(p\) is \(\mathrm{Localization.AtPrime}\,p\)-linearly
  independent.
\end{lemma}

\begin{proof}
  Transport
  \cref{lem:submersive_relation_cotangent_linearIndependent} through the localisation
  map, using that linear independence is preserved by a localised-module map.
\end{proof}

\paragraph{Relative dimension and scheme-to-algebra bridges.}

\begin{lemma}[Existence of a relative dimension for a standard-smooth algebra]\leanok
  \label{lem:exists_standardSmoothOfRelativeDimension_of_standardSmooth}
  \dcref{custom}
  \lean{AlgebraicGeometry.Scheme.exists_isStandardSmoothOfRelativeDimension_of_isStandardSmooth}
  Every standard-smooth \(R\)-algebra \(S\) is standard smooth of some specific
  relative dimension \(n \in \mathbb N\).
\end{lemma}

\begin{proof}
  Take the difference between the number of polynomial variables and
  the number of relations in a submersive presentation of \(S\).
\end{proof}

\begin{lemma}[Submersive presentation of a relative-dimension-$n$ algebra]\leanok
  \label{lem:exists_submersivePresentation_of_standardSmoothReldim}
  \dcref{custom}
  \lean{AlgebraicGeometry.Scheme.exists_submersivePresentation_of_isStandardSmoothOfRelativeDimension}
  \uses{lem:exists_standardSmoothOfRelativeDimension_of_standardSmooth}
  An \(R\)-algebra \(S\) that is standard smooth of relative dimension \(n\) admits an
  explicit submersive presentation \(P\) with \(P.\mathrm{dimension} = n\).
\end{lemma}

\begin{proof}
  By definition, a standard-smooth algebra of relative dimension \(n\)
  is equipped with a submersive presentation of dimension \(n\).
\end{proof}

\begin{lemma}[An affine-open stalk is a localisation]\leanok
  \label{lem:isLocalization_atPrime_stalk_of_affineOpen}
  \dcref{custom}
  \lean{AlgebraicGeometry.Scheme.isLocalization_atPrime_stalk_of_affineOpen}
  For an affine open \(V \ni z\) of a scheme \(X\), the stalk \(\mathcal O_{X,z}\) is
  the localisation of the section ring \(\Gamma(X, V)\) at the prime ideal of \(z\).
\end{lemma}

\begin{proof}
  Identify \(V\) with \(\Spec\Gamma(X,V)\).  Under this identification,
  the stalk at \(z\) is the localisation at the corresponding prime ideal.
\end{proof}

\begin{lemma}[A nonempty open in a one-point space is the whole space]\leanok
  \label{lem:open_eq_top_of_subsingleton}
  \dcref{custom}
  \lean{AlgebraicGeometry.Scheme.open_eq_top_of_subsingleton}
  On a scheme whose underlying space is a subsingleton (at most one point), every
  nonempty open subset equals the whole space.
\end{lemma}

\begin{proof}
  A nonempty open contains the unique point and therefore equals the
  whole underlying space.
\end{proof}

\begin{lemma}[Sections of $\Spec \bar k$ over a nonempty open]\leanok
  \label{lem:gammaSpecField_ringEquiv}
  \dcref{custom}
  \lean{AlgebraicGeometry.Scheme.gammaSpecField_ringEquiv}
  \uses{lem:open_eq_top_of_subsingleton}
  For a field \(\bar k\) and any nonempty open \(U\) of \(\Spec(\bar k)\), the section
  ring \(\Gamma(\Spec \bar k, U)\) is ring-isomorphic to \(\bar k\). This gives the
  algebraically-closed base ring of the standard-smooth presentation a clean
  \(\bar k\)-identification.
\end{lemma}

\begin{proof}
  The spectrum of a field has a single point, so
  \cref{lem:open_eq_top_of_subsingleton} forces \(U = \top\) and the sections collapse
  to the global sections \(\Gamma(\Spec \bar k, \top) \simeq \bar k\).
\end{proof}

\paragraph{Matsumura regular-sequence bridge (Stacks 00NQ).}

\begin{lemma}[Module quotient equals ring quotient]\leanok
  \label{lem:quotSMulTop_quotientRing_linearEquiv}
  \dcref{custom}
  \lean{AlgebraicGeometry.Scheme.quotSMulTop_quotientRing_linearEquiv}
  For \(r\) in a commutative ring \(A\), there is a canonical \((A/(r))\)-linear
  isomorphism between the module quotient \(A/(r \cdot A)\) and the ring quotient
  \(A/(r)\).
\end{lemma}

\begin{proof}
  Both quotients identify elements of \(A\) exactly when their difference
  lies in the principal ideal \((r)\).  The identity on representatives
  therefore induces the required linear isomorphism.
\end{proof}

\begin{lemma}[Ring-quotient cons rule for regular sequences]\leanok
  \label{lem:isRegular_cons_of_quotient_ring}
  \dcref{custom}
  \lean{AlgebraicGeometry.Scheme.isRegular_cons_of_quotient_ring}
  \uses{lem:quotSMulTop_quotientRing_linearEquiv}
  Let \(r \in A\) be such that multiplication by \(r\) is injective on \(A\), and
  suppose the images of a list \(rs\) form a regular sequence over the ring quotient
  \(A/(r)\). Then \(r :: rs\) is a regular sequence over \(A\).
\end{lemma}

\begin{proof}
  Injectivity of multiplication by \(r\) says that \(r\) is a
  non-zero-divisor.  The regularity of the tail in the module quotient
  then gives regularity of the full sequence; use
  \cref{lem:quotSMulTop_quotientRing_linearEquiv} to identify this module
  quotient with \(A/(r)\).
\end{proof}

\begin{lemma}[Cotangent independence descends to a quotient]\leanok
  \label{lem:matsumura_descent_cotangent}
  \dcref{custom}
  \lean{AlgebraicGeometry.Scheme.matsumura_descent_cotangent}
  Let \((A, \mathfrak m)\) be a local ring and \(f_0,\dots,f_m \in \mathfrak m\) with
  \(\kappa(A)\)-linearly independent cotangent classes. Then the tail
  \(f_1,\dots,f_m\), pushed into the quotient \(A' = A/(f_0)\), has
  \(\kappa(A')\)-linearly independent cotangent classes in
  \(\mathfrak m'/\mathfrak m'^2\).
\end{lemma}

\begin{proof}
  The cotangent map \(A \to A'\) is surjective with kernel
  spanned by the class of \(f_0\); a kernel-disjointness argument reduces a
  \(\kappa(A')\)-relation to a \(\kappa(A)\)-relation killed by the full independence
  hypothesis.
\end{proof}

\begin{lemma}[Matsumura: cotangent-independent elements form a regular sequence]\leanok
  \label{lem:matsumura_isRegular_of_linearIndependent_cotangent}
  \dcref{custom}
  \lean{AlgebraicGeometry.Scheme.matsumura_isRegular_of_linearIndependent_cotangent}
  \uses{lem:matsumura_descent_cotangent, lem:isRegular_cons_of_quotient_ring}
  Let \((A, \mathfrak m)\) be a Noetherian regular local ring and \(f_1,\dots,f_n \in
  \mathfrak m\) whose cotangent classes in \(\mathfrak m/\mathfrak m^2\) are
  \(\kappa(A)\)-linearly independent. Then \(f_1,\dots,f_n\) form a regular sequence in
  \(A\) (Matsumura, \emph{Commutative Ring Theory}, Thm.~14.2; Stacks 00NQ).
\end{lemma}

\begin{proof}
  Induct on \(n\). The head \(f_1\) is a non-zero-divisor
  (a regular local ring is a domain), the cotangent independence descends to the
  quotient \(A/(f_1)\) by \cref{lem:matsumura_descent_cotangent}, and the
  ring-quotient cons rule \cref{lem:isRegular_cons_of_quotient_ring} reassembles the
  regular sequence. In the standard-smooth application the linear-independence input is
  supplied by the relations of a submersive presentation via
  \cref{lem:submersive_relation_cotangent_linearIndependent_localized}.
\end{proof}

\begin{lemma}[Smooth over \(\bar k\) \(\Rightarrow\) regular local stalk (Stacks 00TT)]\leanok
  \label{lem:smooth_to_regular_local_ring}
  \dcref{custom}
  \lean{AlgebraicGeometry.Scheme.isRegularLocalRing_stalk_of_smooth}
  \uses{lem:module_free_kaehler_localization, lem:cotangent_kahler_over_field,
        lem:smooth_algebra_krull_dim_formula, lem:stage6_regular_stalk_assembly,
        lem:stalkMap_flat_of_smooth, lem:exists_algebra_standardSmooth_stalk_localization,
        lem:finrank_cotangentSpace_of_bijective_residue, lem:gammaSpecField_ringEquiv,
        lem:exists_standardSmoothOfRelativeDimension_of_standardSmooth}
  % SOURCE: \cite{stacks-project}, tag 00TT --- Jacobian criterion
  % (smooth-implies-regular direction) (read from references/stacks-algebra.tex,
  % L38593--38611; tag verified against the Stacks website tag/00TT page,
  % "Lemma 10.140.3" in section 10.140 "Smooth algebras over fields").
  % SOURCE QUOTE: "Let $k$ be any field. Let $S$ be a finite type $k$-algebra.
  % Let $X = \Spec(S)$. Let $\mathfrak q \subset S$ be a prime corresponding to
  % $x \in X$. The following are equivalent:
  % \begin{enumerate}
  % \item The $k$-algebra $S$ is smooth at $\mathfrak q$ over $k$.
  % \item We have
  % $\dim_{\kappa(\mathfrak q)} \Omega_{S/k} \otimes_S \kappa(\mathfrak q)
  % \leq \dim_x X$.
  % \item We have
  % $\dim_{\kappa(\mathfrak q)} \Omega_{S/k} \otimes_S \kappa(\mathfrak q)
  % = \dim_x X$.
  % \end{enumerate}
  % Moreover, in this case the local ring $S_{\mathfrak q}$ is regular."
  \textit{Source: \cite{stacks-project}, tag 00TT.}
  Let \(X\) be a variety over an algebraically closed field \(\bar k\) (geometrically
  integral, separated, locally of finite type) and assume the structure morphism
  \(X \to \Spec(\bar k)\) is \emph{smooth}. Then for every point \(z \in X\) the
  stalk \(\mathcal O_{X, z}\) is a regular local ring.

  The algebraic-closure hypothesis ensures that every closed point is
  \(\bar k\)-rational. Regularity at an arbitrary point may then also be obtained by
  choosing a closed specialization and localizing its regular local ring. The
  arbitrary-field form follows directly from the general smooth-implies-regular
  criterion of tag \texttt{00TT}.
\end{lemma}

\begin{proof}
  \textit{Source: \cite{stacks-project}, tag 00TT.} The cited lemma packages the
  Jacobian-criterion equivalence and concludes ``in this case the local ring
  \(S_{\mathfrak q}\) is regular''. Translated to the scheme level: a smooth
  finite-type morphism \(X \to \Spec(\bar k)\) restricted to an affine chart
  \(\Spec(S) \subseteq X\) gives an affine algebra \(\bar k \to S\) that is
  smooth at every prime \(\mathfrak q \subset S\); tag 00TT then yields
  regularity of the localisation \(S_{\mathfrak q} = \mathcal O_{X, z}\) at every
  point \(z \in X\) lying over \(\mathfrak q\).
\end{proof}

\begin{lemma}[Smooth at a codim-1 point $\Rightarrow$ principal nonzero maximal ideal]\leanok
  \label{lem:smooth_codim_one_maximalIdeal_principal}
  \dcref{custom}
  \lean{AlgebraicGeometry.Scheme.smooth_codim_one_maximalIdeal_isPrincipal_and_ne_bot}
  \uses{lem:smooth_to_regular_local_ring, thm:ringKrullDim_stalk_eq_coheight}
  Let \(X\) be a smooth, geometrically-irreducible, separated, integral variety over an
  algebraically closed field \(\bar k\), and let \(z \in X\) be a point with
  \(\mathrm{coheight}\,z = 1\) (a codimension-one point). Then the maximal ideal of the
  stalk \(\mathcal O_{X,z}\) is principal and nonzero.
\end{lemma}

\begin{proof}
  Smoothness gives a regular local stalk via
  \cref{lem:smooth_to_regular_local_ring}; the coheight bridge
  \cref{thm:ringKrullDim_stalk_eq_coheight} gives Krull dimension \(1\); the
  cotangent-space characterisation of regularity then forces
  \(\dim_\kappa \mathfrak m/\mathfrak m^2 = 1\), i.e.\ a principal maximal ideal, which
  is nonzero since the stalk is not a field. This is the geometric core consumed by
  \cref{lem:smooth_codim_one_dvr}.
\end{proof}

\begin{lemma}[Smooth \(\Rightarrow\) DVR at every codim-\(1\) point]\leanok
  \label{lem:smooth_codim_one_dvr}
  \dcref{custom}
  \lean{AlgebraicGeometry.Scheme.localRing_dvr_of_codim_one}
  \uses{lem:smooth_to_regular_local_ring, thm:ringKrullDim_stalk_eq_coheight,
        lem:smooth_codim_one_maximalIdeal_principal}
  % SOURCE: \cite{hartshorne-ag}, II.6, p.~130 (the definition ``regular in codimension one'',
  % together with the immediately preceding remark that on a nonsingular variety every
  % local ring is regular) (read from references/hartshorne-algebraic-geometry.pdf,
  % PDF page 147).
  % SOURCE QUOTE: "\textbf{Definition.} We say a scheme \(X\) is \emph{regular in
  % codimension one} (or sometimes \emph{nonsingular in codimension one}) if every
  % local ring \(\mathcal O_x\) of \(X\) of dimension one is regular."
  % SOURCE QUOTE: "The most important examples of such schemes are nonsingular
  % varieties over a field (I, \S5) and noetherian normal schemes. On a nonsingular
  % variety the local ring of every closed point is regular (I, 5.1), hence all the
  % local rings are regular, since they are localizations of the local rings of
  % closed points."
  % SOURCE QUOTE: "If \(Y\) is a prime divisor on \(X\), let \(\eta \in Y\) be its generic
  % point. Then the local ring \(\mathcal O_{\eta, X}\) of \(X\) at \(\eta\) is a discrete
  % valuation ring with quotient field \(K\), the function field of \(X\)."
  \textit{Source: Hartshorne, II.6, p.~130 (definition of regular in codimension one;
  the DVR consequence at a codim-1 point).}
  Let \(X\) be a nonsingular variety over \(\bar k\) and let \(\eta \in X\) be a point with
  \(\mathrm{codim}_X \overline{\{\eta\}} = 1\). Then the local ring \(\mathcal O_{X, \eta}\) is a
  discrete valuation ring with fraction field the function field \(K(X)\) of \(X\).
\end{lemma}

\begin{proof}

  By \cref{lem:smooth_to_regular_local_ring} (Stacks tag \texttt{00TT}, the
  Jacobian-criterion direction), the stalk \(\mathcal O_{X, x}\) is regular for
  every point \(x \in X\) --- in particular at every closed point (the case
  Hartshorne~I.5.1 states classically). Localising a regular local ring at a
  prime ideal yields a regular local ring; in particular, for \(\eta \in X\) with
  \(\mathrm{codim}_X \overline{\{\eta\}} = 1\), the local ring
  \(\mathcal O_{X, \eta}\) is regular of Krull dimension \(1\). A regular local ring of Krull dimension \(1\) is exactly a
  discrete valuation ring (Atiyah--Macdonald, Proposition 9.2; or Stacks tag
  \texttt{00PD}). Its fraction field is the function field \(K(X)\) of \(X\) by
  irreducibility of \(X\) (the unique generic point of \(X\) lies above \(\eta\), and the
  stalk at the generic point of an integral scheme is the function field).
\end{proof}

\section{Milne Theorem 3.1: extension across codimension at least two}
\label{sec:milne_thm31}

Before the extension theorem itself we record the function-field functoriality
\(K(Y) \to K(X)\) attached to a dominant rational map, which the codim-\(1\) step of
the extension proof uses to view a representative as a morphism \(\Spec K(X) \to Y\).

\begin{lemma}[Image of the closed point of \(\Spec K(X)\)]\leanok
  \label{lem:fromFunctionField_base_eq_genericPoint}
  \dcref{custom}
  \lean{AlgebraicGeometry.Scheme.RationalMap.fromFunctionField_base_eq_genericPoint}
  Let \(X\) be an integral scheme and \(Y\) an irreducible scheme, and let
  \(f \colon X \dashrightarrow Y\) be a dominant rational map. The induced morphism
  \(\Spec K(X) \to Y\) (the restriction of \(f\) to the generic point of \(X\))
  sends the closed point of \(\Spec K(X)\) to the generic point of \(Y\).
\end{lemma}

\begin{proof}\leanok
  Choose a representative \((U, \varphi_U)\) of \(f\); its underlying morphism
  \(\varphi_U \colon U \to Y\) is dominant. The morphism \(\Spec K(X) \to Y\) factors as
  \(\Spec \mathcal O_{X,\eta_X} \to U \xrightarrow{\varphi_U} Y\), where the first map
  sends the closed point to the generic point \(\eta_U\) of \(U\) (checked after the open
  immersion \(U \hookrightarrow X\), which identifies \(\eta_U\) with \(\eta_X\)). A
  dominant morphism carries the generic point of its irreducible source to a generic
  point of \(Y\); as \(Y\) is a \(T_0\) space its generic point is unique, so the image
  is \(\eta_Y\).
\end{proof}

\begin{definition}\leanok
  \label{def:functionFieldPullback}
  \dcref{custom}
  \lean{AlgebraicGeometry.Scheme.RationalMap.functionFieldPullback}
  \uses{lem:fromFunctionField_base_eq_genericPoint}
  For a dominant rational map \(f \colon X \dashrightarrow Y\) from an integral scheme
  \(X\) to an irreducible scheme \(Y\), the \textbf{function-field pullback} is the ring
  homomorphism \(K(Y) \to K(X)\) obtained from the germ pullback
  \(\mathcal O_{Y, f(\eta_X)} \to K(X)\) of \(f\) at the generic point, identified along
  \cref{lem:fromFunctionField_base_eq_genericPoint} with the stalk
  \(\mathcal O_{Y,\eta_Y} = K(Y)\).
\end{definition}

The difference map of Milne's Lemma~3.3 (\cref{sec:milne_lem33} below) is built by
precomposing \(f\) with the two projections \(X \times X \to X\) and pairing the results.
The precomposition operation is the following; it is the left-hand analogue of the
right-composition of a rational map with a morphism, \(f \mapsto f \circ g\).

\begin{definition}\leanok
  \label{def:rationalMap_precomp}
  \dcref{custom}
  \lean{AlgebraicGeometry.Scheme.RationalMap.precomp}
  Let \(p \colon W \to X\) be a morphism of schemes whose underlying continuous map is
  \emph{open}, and let \(f \colon X \dashrightarrow Y\) be a rational map. The
  \textbf{precomposition} \(f \circ p \colon W \dashrightarrow Y\) is the rational map
  represented, for any partial-map representative \((V, \varphi_V)\) of \(f\), by
  \[
    \bigl(p^{-1}(V),\; \; p^{-1}(V) \xrightarrow{\;p|_{p^{-1}(V)}\;} V
    \xrightarrow{\;\varphi_V\;} Y\bigr).
  \]
  The domain \(p^{-1}(V)\) is dense because the preimage of a dense open set under an open
  continuous map is dense; the same fact makes the construction independent of the chosen
  representative, so it descends to a well-defined rational map.
\end{definition}

\begin{lemma}\leanok
  \label{lem:rationalMap_precomp_compHom}
  \dcref{custom}
  \lean{AlgebraicGeometry.Scheme.RationalMap.precomp_compHom}
  \uses{def:rationalMap_precomp}
  Precomposition on the left commutes with right-composition by a morphism: for
  \(p \colon W \to X\) open, \(f \colon X \dashrightarrow Y\), and \(g \colon Y \to Z\),
  \[
    (f \circ p) \circ g \;=\; (f \circ g) \circ p
    \qquad\text{in}\qquad W \dashrightarrow Z .
  \]
  Furthermore, if \(p\) and \(f\) are defined over a base scheme \(S\), then so is
  \(f \circ p\).
\end{lemma}
\begin{proof}\leanok
  Both sides act on a chosen partial-map representative \((V, \varphi_V)\) of \(f\) by
  restriction along \(p\) and composition with \(g\); the resulting equality is the
  associativity of morphism composition,
  \((p|_{p^{-1}(V)} \circ \varphi_V) \circ g = p|_{p^{-1}(V)} \circ (\varphi_V \circ g)\).
  For the over-\(S\) claim, \(p|_{p^{-1}(V)}\) is an \(S\)-morphism because it commutes with
  the structure morphisms via the restriction identity
  \(p|_{p^{-1}(V)} \circ (V \hookrightarrow X) = (p^{-1}(V) \hookrightarrow W) \circ p\),
  and a composite of \(S\)-morphisms is an \(S\)-morphism.
\end{proof}

\begin{lemma}\leanok
  \label{lem:rationalMap_precomp_hom}
  \dcref{custom}
  \lean{AlgebraicGeometry.Scheme.RationalMap.precomp_hom_toRationalMap}
  \uses{def:rationalMap_precomp}
  Precomposing the rational map of an honest morphism \(g \colon X \to Y\) with an open
  morphism \(p \colon W \to X\) is the rational map of the composite:
  \[
    (g \text{ as a rational map}) \circ p \;=\; (p \circ g) \text{ as a rational map} .
  \]
\end{lemma}
\begin{proof}\leanok
  Both sides are everywhere-defined; on the whole of \(W\) the left-hand side is
  \(p \circ g\) by the restriction identity \(p|_{W} \circ (X \hookrightarrow X)_{\mathrm{id}}
  = p\), i.e.\ the two partial maps have the same (dense, in fact total) domain and equal
  underlying morphism.
\end{proof}

To pair the two precompositions \(f \circ \mathrm{pr}_1\) and \(f \circ \mathrm{pr}_2\) into
a single rational map \(X \times X \dashrightarrow G \times_{\bar k} G\), we use the
function-field correspondence: for an integral scheme \(X\) over \(S\) and a scheme \(Y\)
locally of finite type over \(S\), a rational map \(X \dashrightarrow Y\) over \(S\) is the
same datum as an \(S\)-morphism \(\operatorname{Spec} K(X) \to Y\) (restriction to the generic
point, spread out again by local finite type). Morphisms out of the single scheme
\(\operatorname{Spec} K(X)\) pair freely through the universal property of the fibre product.

\begin{definition}\leanok
  \label{def:rationalMap_prod}
  \dcref{custom}
  \lean{AlgebraicGeometry.Scheme.RationalMap.prod}
  Let \(s_X \colon X \to S\), \(s_Y \colon Y \to S\), \(s_Z \colon Z \to S\) be structure
  morphisms with \(X\) integral and \(s_Y, s_Z\) locally of finite type, and let
  \(a \colon X \dashrightarrow Y\), \(b \colon X \dashrightarrow Z\) be rational maps over
  \(S\) (i.e.\ \(a\) restricted to the generic point commutes with the structure morphisms,
  and likewise for \(b\)). The \textbf{pairing} \(a \times_S b \colon X \dashrightarrow
  Y \times_S Z\) is the rational map corresponding, under the function-field correspondence,
  to the morphism
  \[
    \operatorname{Spec} K(X)
    \xrightarrow{\;\langle a_\eta,\, b_\eta \rangle\;}
    Y \times_S Z
  \]
  obtained from the generic-point morphisms \(a_\eta \colon \operatorname{Spec} K(X) \to Y\)
  and \(b_\eta \colon \operatorname{Spec} K(X) \to Z\) by the universal property of the fibre
  product (their composites with the structure morphisms both equal
  \(\operatorname{Spec} K(X) \to \operatorname{Spec}\mathcal O_{X,\eta} \to X \to S\)).
\end{definition}

\begin{lemma}\leanok
  \label{lem:rationalMap_prod_proj}
  \dcref{custom}
  \lean{AlgebraicGeometry.Scheme.RationalMap.prod_compHom_fst,
        AlgebraicGeometry.Scheme.RationalMap.prod_compHom_snd,
        AlgebraicGeometry.Scheme.RationalMap.prod_compHom_over}
  \uses{def:rationalMap_prod}
  The two projections of the pairing recover the factors, and the pairing is again over
  \(S\):
  \[
    (a \times_S b) \circ \mathrm{pr}_1 = a,
    \qquad
    (a \times_S b) \circ \mathrm{pr}_2 = b,
    \qquad
    (a \times_S b) \circ \bigl(\mathrm{pr}_1 \circ s_Y\bigr) = s_X,
  \]
  where \(\mathrm{pr}_1 \colon Y \times_S Z \to Y\), \(\mathrm{pr}_2 \colon Y \times_S Z \to Z\)
  are the projections (the last equation being the statement that \(a \times_S b\) is over
  \(S\)).
\end{lemma}
\begin{proof}\leanok
  Right-composition with a morphism is functorial on generic-point morphisms:
  \((\varphi \circ g)_\eta = \varphi_\eta \circ g\) for any rational map \(\varphi\) and
  morphism \(g\). Hence the generic-point morphism of \((a \times_S b) \circ \mathrm{pr}_1\) is
  \(\langle a_\eta, b_\eta\rangle \circ \mathrm{pr}_1 = a_\eta\), which is the generic-point
  morphism of \(a\); since a rational map out of an integral scheme is determined by its
  generic-point morphism, \((a \times_S b) \circ \mathrm{pr}_1 = a\). The second identity is
  symmetric. The over-\(S\) identity follows by combining the first identity with
  associativity of right-composition, \((a \times_S b) \circ (\mathrm{pr}_1 \circ s_Y) =
  \bigl((a \times_S b) \circ \mathrm{pr}_1\bigr) \circ s_Y = a \circ s_Y = s_X\).
\end{proof}

The two precompositions and their pairing are assembled into Milne's difference map by
right-composing with the \emph{difference morphism} of the group object \(G\), the honest
group law \((g,h) \mapsto g \cdot h^{-1}\).

\begin{definition}\leanok
  \label{def:grpObj_diff}
  \dcref{custom}
  \lean{CategoryTheory.GrpObj.diff, AlgebraicGeometry.grpObjDiffLeft}
  Let \(G\) be a group object in a cartesian-monoidal category. The \textbf{difference
  morphism} \(\operatorname{diff}_G \colon G \otimes G \to G\) is
  \(\operatorname{fst} / \operatorname{snd}\), the division in the group structure on the
  hom-set \(\operatorname{Hom}(G \otimes G, G)\); pointwise it is
  \((g,h) \mapsto g \cdot h^{-1}\), i.e.\ the group law composed with inversion on the second
  factor. Equivalently, it is built from the group structure on the hom-set. When \(G\)
  lives over \(\operatorname{Spec}\bar k\), its underlying scheme morphism
  \(\operatorname{diff}_G^{\,\mathrm{left}} \colon G \times_{\bar k} G \to G\) is Milne's
  \(m \circ (\mathrm{id} \times \mathrm{inv})\).
\end{definition}

\begin{lemma}\leanok
  \label{lem:grpObjDiffLeft_comp_hom}
  \dcref{custom}
  \lean{AlgebraicGeometry.grpObjDiffLeft_comp_hom}
  \uses{def:grpObj_diff}
  The difference morphism is over the base: \(\operatorname{diff}_G^{\,\mathrm{left}} \circ
  (G \to \operatorname{Spec}\bar k) = \mathrm{pr}_1 \circ (G \to \operatorname{Spec}\bar k)\).
\end{lemma}
\begin{proof}\leanok
  This is the triangle identity \(\varphi_{\mathrm{left}} \circ g_{\mathrm{hom}} =
  f_{\mathrm{hom}}\) of the over-category morphism \(\operatorname{diff}_G\), together with the
  identification of the structure morphism of \(G \times_{\bar k} G\) as
  \(\mathrm{pr}_1 \circ (G \to \operatorname{Spec}\bar k)\).
\end{proof}

\begin{definition}\leanok
  \label{def:rationalMap_difference}
  \dcref{custom}
  \lean{AlgebraicGeometry.Scheme.RationalMap.differenceRationalMap}
  \uses{def:rationalMap_precomp, def:rationalMap_prod, def:grpObj_diff}
  Let \(X\) be smooth over an algebraically closed field \(\bar k\) with \(X \times_{\bar k} X\)
  integral, let \(G\) be a group variety over \(\bar k\), and let
  \(f \colon X \dashrightarrow G\) be a rational map over \(\bar k\). Milne's
  \textbf{difference rational map}
  \[
    \Phi \;=\; \bigl((f \circ \mathrm{pr}_1) \times_{\bar k} (f \circ \mathrm{pr}_2)\bigr)
    \circ \operatorname{diff}_G^{\,\mathrm{left}}
    \colon\; X \times_{\bar k} X \;\dashrightarrow\; G
  \]
  is \((x,y) \mapsto f(x) \cdot f(y)^{-1}\). The projections \(\mathrm{pr}_i\) are open (being
  base changes of the smooth structure morphism), so the precompositions
  \(f \circ \mathrm{pr}_i\) are defined; both are over \(\bar k\)
  (\cref{lem:rationalMap_precomp_compHom}, \cref{lem:rationalMap_precomp_hom}) with common
  structure composite \(\mathrm{pr}_1 \circ (X \to \operatorname{Spec}\bar k)\), since
  \(\mathrm{pr}_1\) and \(\mathrm{pr}_2\) agree after \(X \to \operatorname{Spec}\bar k\).
\end{definition}

\begin{lemma}\leanok
  \label{lem:rationalMap_difference_over}
  \dcref{custom}
  \lean{AlgebraicGeometry.Scheme.RationalMap.differenceRationalMap_compHom_over}
  \uses{def:rationalMap_difference, lem:rationalMap_prod_proj, lem:grpObjDiffLeft_comp_hom}
  The difference map is a \(\bar k\)-rational map: \(\Phi \circ (G \to \operatorname{Spec}\bar k)
  = \mathrm{pr}_1 \circ (X \to \operatorname{Spec}\bar k)\), the structure morphism of
  \(X \times_{\bar k} X\).
\end{lemma}
\begin{proof}\leanok
  By associativity of right-composition, \(\Phi \circ (G \to \operatorname{Spec}\bar k) =
  \bigl((f \circ \mathrm{pr}_1) \times (f \circ \mathrm{pr}_2)\bigr) \circ
  \bigl(\operatorname{diff}_G^{\,\mathrm{left}} \circ (G \to \operatorname{Spec}\bar k)\bigr)\);
  by \cref{lem:grpObjDiffLeft_comp_hom} the inner composite is
  \(\mathrm{pr}_1 \circ (G \to \operatorname{Spec}\bar k)\), and the over-\(S\) identity of the
  pairing (\cref{lem:rationalMap_prod_proj}) evaluates the result to the structure morphism of
  \(X \times_{\bar k} X\).
\end{proof}

The following elementary fact is the workhorse for Milne's ``\(\Phi\) is defined at
\((x,x)\) if \(f\) is defined at \(x\)'' step: post-composing a rational map with an
honest morphism never destroys definedness.

\begin{lemma}\leanok
  \label{lem:compHom_domain_monotone}
  \dcref{custom}
  \lean{AlgebraicGeometry.Scheme.RationalMap.le_domain_compHom}
  Right-composition of a rational map \(F \colon X \dashrightarrow Y\) with a morphism
  \(g \colon Y \to Z\) enlarges the domain of definition:
  \(\mathrm{Dom}(F) \subseteq \mathrm{Dom}(F \circ g)\).
\end{lemma}
\begin{proof}\leanok
  A partial-map representative \((V, F_V)\) of \(F\) yields the representative
  \((V, F_V \circ g)\) of \(F \circ g\) with the same domain \(V\).
\end{proof}

\begin{lemma}\leanok
  \label{lem:difference_map_domain_lower_bound}
  \dcref{custom}
  \lean{AlgebraicGeometry.Scheme.RationalMap.le_domain_differenceRationalMap}
  \uses{def:rationalMap_difference, def:rationalMap_precomp, def:rationalMap_prod,
        lem:compHom_domain_monotone}
  \textbf{(Milne Lemma~3.3, Sub-step~1 domain bound / Sub-step~2 easy direction.)} The
  difference rational map \(\Phi\) is defined on \(\mathrm{Dom}(f) \times_{\bar k}
  \mathrm{Dom}(f)\): every point of \(X \times_{\bar k} X\) projecting into
  \(\mathrm{Dom}(f)\) under both projections lies in \(\mathrm{Dom}(\Phi)\). In
  particular \(\Phi\) is defined at the diagonal point \((x, x)\) whenever \(f\) is
  defined at \(x\).
\end{lemma}
\begin{proof}\leanok
  Right-composition with \(\operatorname{diff}_G\) only enlarges the domain of
  definition (\cref{lem:compHom_domain_monotone}), so it suffices to show the pairing
  \((f \circ \mathrm{pr}_1) \times (f \circ \mathrm{pr}_2)\) is defined on
  \(\mathrm{Dom}(f) \times_{\bar k} \mathrm{Dom}(f)\). The maximal representative
  \((U, \varphi_U)\) of \(f\) (on its full domain of definition) is defined on \(U =
  \mathrm{Dom}(f)\) and over \(\bar k\); pairing \(\varphi_U \circ \mathrm{pr}_i\) into
  \(G \times_{\bar k} G\) over the common domain
  \(\mathrm{pr}_1^{-1}U \cap \mathrm{pr}_2^{-1}U\) gives Milne's explicit morphism
  \(U \times_{\bar k} U \to G \times_{\bar k} G\), a partial map whose induced
  function-field morphism on the integral source \(X \times_{\bar k} X\) matches that of
  the pairing (\cref{def:rationalMap_prod}) factor-by-factor. Hence it represents the
  pairing, and its domain \(\mathrm{pr}_1^{-1}U \cap \mathrm{pr}_2^{-1}U\) lies in the
  pairing's domain of definition.
\end{proof}

The converse direction of Milne's slice argument rests on the reconstruction formula
\(f(x) = \Phi(x,u) \cdot f(u)\). The next four lemmas establish its group-theoretic
content and the resulting domain bound.  Together with openness of
\(\mathrm{Dom}(\Phi)\), irreducibility of the fibre \(X_{\kappa(x)}\), and descent of
definedness along the smooth projection, this proves that definedness of \(\Phi\) at
\((x,x)\) forces definedness of \(f\) at \(x\).

\begin{lemma}\leanok
  \label{lem:rationalMap_precomp_domain_bound}
  \dcref{custom}
  \lean{AlgebraicGeometry.Scheme.RationalMap.le_domain_precomp}
  \uses{def:rationalMap_precomp}
  For a rational map \(f \colon X \dashrightarrow Y\) and an open morphism
  \(p \colon W \to X\), the precomposition is defined on the preimage of the domain:
  \(p^{-1}(\mathrm{Dom}(f)) \subseteq \mathrm{Dom}(f \circ p)\).
\end{lemma}
\begin{proof}\leanok
  If \(p(w) \in \mathrm{Dom}(f)\), choose a partial-map representative \((V, f_V)\) of
  \(f\) with \(p(w) \in V\). Then \((p^{-1}V, f_V \circ (p \mid_V))\) is a representative
  of \(f \circ p\) defined at \(w\), so \(w \in \mathrm{Dom}(f \circ p)\).
\end{proof}

\begin{lemma}\leanok
  \label{lem:rationalMap_prod_domain_bound}
  \dcref{custom}
  \lean{AlgebraicGeometry.Scheme.RationalMap.pairPartialMap,
        AlgebraicGeometry.Scheme.RationalMap.le_domain_prod}
  \uses{def:rationalMap_prod}
  Let \(X\) be integral, let \(S\), \(Y\), \(Z\) be separated schemes with structure
  morphisms \(s_Y \colon Y \to S\), \(s_Z \colon Z \to S\) locally of finite type, and let
  \(a \colon X \dashrightarrow Y\), \(b \colon X \dashrightarrow Z\) be rational maps over
  \(S\). Then the pairing is defined on the intersection of the two domains:
  \(\mathrm{Dom}(a) \cap \mathrm{Dom}(b) \subseteq \mathrm{Dom}(a \times_S b)\).
\end{lemma}
\begin{proof}\leanok
  The maximal representatives of \(a\) and \(b\) are defined on \(\mathrm{Dom}(a)\),
  resp.\ \(\mathrm{Dom}(b)\), and are \(S\)-morphisms; over the common domain
  \(\mathrm{Dom}(a) \cap \mathrm{Dom}(b)\) they agree after the structure morphisms to
  \(S\), hence pair through the fibre product into a partial map defined on
  \(\mathrm{Dom}(a) \cap \mathrm{Dom}(b)\). Its induced function-field morphism restricts,
  on each fibre-product factor, to \(a\)'s, resp.\ \(b\)'s; so it represents \(a \times_S
  b\), and its domain lies in the domain of definition of the pairing.
\end{proof}

\begin{lemma}\leanok
  \label{lem:grpObj_reconstruction}
  \dcref{custom}
  \lean{CategoryTheory.GrpObj.lift_diff_lift_mul}
  \uses{def:grpObj_diff}
  Let \(G\) be a group object in a cartesian-monoidal category and \(a, b \colon T \to G\)
  two \(T\)-points. Then
  \(\bigl\langle \langle a, b\rangle \circ \operatorname{diff}_G,\ b \bigr\rangle \circ m
  = a\), i.e.\ \((a \cdot b^{-1}) \cdot b = a\).
\end{lemma}
\begin{proof}\leanok
  Precomposition by \(\langle a, b\rangle\) is a group homomorphism on \(T\)-points, so
  \(\langle a, b\rangle \circ \operatorname{diff}_G = \langle a, b\rangle \circ
  (\operatorname{fst} / \operatorname{snd}) = a / b\). Then \(\langle a/b, b\rangle \circ m
  = (a/b) \cdot b = a\) by the group cancellation law.
\end{proof}

\begin{lemma}\leanok
  \label{lem:scheme_reconstruction}
  \dcref{custom}
  \lean{AlgebraicGeometry.grpObjMulLeft,
        AlgebraicGeometry.pullback_lift_diff_lift_mul}
  \uses{lem:grpObj_reconstruction, def:grpObj_diff}
  Let \(G\) be a group variety over \(\bar k\) and \(A, B \colon T \to G\) two
  \(\bar k\)-morphisms sharing a structure map (\(A \circ (G \to \operatorname{Spec}\bar k)
  = B \circ (G \to \operatorname{Spec}\bar k)\)). Then the scheme-level reconstruction holds:
  \[
    m^{\mathrm{left}} \circ \bigl\langle
      \operatorname{diff}_G^{\,\mathrm{left}} \circ \langle A, B\rangle,\ B
    \bigr\rangle \;=\; A .
  \]
\end{lemma}
\begin{proof}\leanok
  The forgetful functor \((-)^{\mathrm{left}} \colon \mathbf{Sch}/\bar k \to \mathbf{Sch}\)
  sends the cartesian lift, first and second projections, and composition to
  \(\mathrm{pullback}\)-lift, \(\mathrm{pr}_1, \mathrm{pr}_2\) and composition. Lifting \(A,
  B\) to \(\bar k\)-morphisms \(T \to G\) in the over-category and applying
  \((-)^{\mathrm{left}}\) to \cref{lem:grpObj_reconstruction} yields the claim.
\end{proof}

\begin{lemma}\leanok
  \label{lem:rationalMap_difference_reconstruct}
  \dcref{custom}
  \lean{AlgebraicGeometry.Scheme.RationalMap.reconstruct_precomp_fst}
  \uses{def:rationalMap_difference, def:rationalMap_prod, lem:scheme_reconstruction}
  The reconstruction \(f(x) = \Phi(x,y) \cdot f(y)\) holds as an identity of rational maps
  on \(X \times_{\bar k} X\):
  \[
    f \circ \mathrm{pr}_1 \;=\; m^{\mathrm{left}} \circ
    \bigl(\Phi \times_{\bar k} (f \circ \mathrm{pr}_2)\bigr).
  \]
\end{lemma}
\begin{proof}\leanok
  Both sides are \(\bar k\)-rational maps out of the integral scheme \(X \times_{\bar k}
  X\), so it suffices to match their function-field morphisms \(\operatorname{Spec}
  K(X \times X) \to G\). The left side is \(A = (f \circ \mathrm{pr}_1)^{\mathrm{ff}}\); the
  right side unfolds, via the function-field formula for the pairing and difference map, to
  \(m^{\mathrm{left}} \circ \langle \operatorname{diff}_G^{\,\mathrm{left}} \circ \langle A,
  B\rangle,\ B\rangle\) with \(B = (f \circ \mathrm{pr}_2)^{\mathrm{ff}}\). These agree by
  \cref{lem:scheme_reconstruction}.
\end{proof}

\begin{lemma}\leanok
  \label{lem:difference_map_domain_reconstruction}
  \dcref{custom}
  \lean{AlgebraicGeometry.Scheme.RationalMap.le_domain_precomp_fst_of_difference}
  \uses{lem:rationalMap_difference_reconstruct, lem:rationalMap_precomp_domain_bound,
        lem:rationalMap_prod_domain_bound, lem:compHom_domain_monotone}
  \textbf{(Milne Lemma~3.3, Sub-step~2 hard direction, algebraic content.)}
  \(f \circ \mathrm{pr}_1\) is defined wherever \(\Phi\) is defined and \(f \circ
  \mathrm{pr}_2\) is defined:
  \[
    \mathrm{Dom}(\Phi) \cap \mathrm{pr}_2^{-1}(\mathrm{Dom}(f)) \;\subseteq\;
    \mathrm{Dom}(f \circ \mathrm{pr}_1) .
  \]
  This is the group-theoretic half of Milne's converse: at a point where \(\Phi\) is
  defined and \(f\) is defined on the second coordinate, the reconstruction
  \(f(x) = \Phi(x,u) \cdot f(u)\) exhibits \(f \circ \mathrm{pr}_1\) as regular.
\end{lemma}
\begin{proof}\leanok
  Rewriting \(f \circ \mathrm{pr}_1\) by
  \cref{lem:rationalMap_difference_reconstruct}, it suffices to bound
  \(\mathrm{Dom}\bigl(m^{\mathrm{left}} \circ (\Phi \times (f \circ \mathrm{pr}_2))\bigr)\).
  Right-composition only enlarges the domain (\cref{lem:compHom_domain_monotone}), and the
  pairing is defined on \(\mathrm{Dom}(\Phi) \cap \mathrm{Dom}(f \circ \mathrm{pr}_2)\)
  (\cref{lem:rationalMap_prod_domain_bound}), which contains \(\mathrm{Dom}(\Phi) \cap
  \mathrm{pr}_2^{-1}(\mathrm{Dom}(f))\) by \cref{lem:rationalMap_precomp_domain_bound}.
\end{proof}

The first input to Milne's Theorem~3.2 is the following: a rational map from a normal
variety to a \emph{complete} variety automatically extends across every codim-\(1\)
point.

\begin{theorem}[Milne Theorem~3.1: codim-\(\geq 2\) extension to a complete target]\leanok
  \label{thm:codim_one_extension}
  \dcref{custom}
  \lean{AlgebraicGeometry.Scheme.RationalMap.indeterminacy_codimGe2_of_smooth_of_complete}
  \uses{def:indeterminacy_locus, def:codim_one_indeterminacy,
        lem:smooth_codim_one_dvr}
  % SOURCE: \cite{milne-av}, Theorem~3.1, \S I.3, p.~16
  % (read from references/abelian-varieties.pdf, PDF page 22).
  % SOURCE QUOTE: "\textsc{Theorem 3.1.} \emph{A rational map $\varphi \colon V
  % \dashrightarrow W\( from a normal variety \)V\( to a complete variety \)W$ is defined
  % on an open subset \(U\) of \(V\) whose complement \(V \smallsetminus U\) has codimension
  % \(\geq 2\).}"
  \textit{Source: Milne, \emph{Abelian Varieties}, Theorem~3.1, \S I.3, p.~16.}
  Let \(X\) be a nonsingular variety over \(\bar k\), let \(Y\) be a complete variety over
  \(\bar k\), and let \(f \colon X \dashrightarrow Y\) be a rational map. Then the
  indeterminacy locus \(Z(f)\) (\cref{def:indeterminacy_locus}) has codimension
  \(\geq 2\) in \(X\); equivalently, \(f\) is codim-\(1\)-indeterminacy-free
  (\cref{def:codim_one_indeterminacy}).
\end{theorem}

% SOURCE QUOTE PROOF: "\textsc{Proof.} Assume first that \(V\) is a curve. Thus we are
% given a nonsingular curve \(V\) and a regular map \(\varphi \colon U \to W\) from an
% open subset of \(V\) which we want to extend to \(V\). Consider the maps
% [diagram: \(U \hookrightarrow V\), $U \xrightarrow{u \mapsto (u, \varphi(u))} Z
% \subseteq V \times W\(, \)U \xrightarrow{\varphi} W$, with projections from
% \(V \times W\) to \(V\) and to \(W\)]
% Let \(U'\) be the image of \(U\) in \(V \times W\), and let \(Z\) be its closure. The image
% of \(Z\) in \(V\) is closed (because \(W\) is complete), and contains \(U\) (the composite
% \(U \to V\) is the given inclusion), and so \(Z\) maps onto \(V\). The maps $U \to U'
% \to U\( are isomorphisms. Therefore, \)Z \to V$ is a surjective map from a complete
% curve onto a nonsingular complete curve that is an isomorphism on open subsets.
% Such a map must be an isomorphism (complete nonsingular curves are determined by
% their function fields). The restriction of the projection map \(V \times W \to W\)
% to \(Z\ (\simeq V)\) is the extension of \(\varphi\) to \(V\) we are seeking.
%   The general case can be reduced to the one-dimensional case (using schemes). Let
% \(U\) be the largest subset on which \(\varphi\) is defined, and suppose that $V
% \smallsetminus U\( has codimension 1. Then there is a prime divisor \)Z\( in \)V
% \smallsetminus U\(. Because \)V$ is normal, its associated local ring is a discrete
% valuation ring \(\mathcal O_Z\) with field of fractions \(k(V)\). The map \(\varphi\)
% defines a morphism of schemes \(\operatorname{Spec}(k(V)) \to W\), which the
% valuative criterion of properness (Hartshorne 1977, II 4.7) shows extends to a
% morphism \(\operatorname{Spec}(\mathcal O_Z) \to W\). This implies that \(\varphi\)
% has a representative defined on an open subset that meets \(Z\) in a nonempty set,
% which is a contradiction."
\begin{proof}\leanok
  Fix an indeterminacy point \(z \in Z(f)\); we must show its coheight is \(\geq 2\).
  Suppose not, so \(\mathrm{coheight}_X(z) \leq 1\), and argue by cases.

  \emph{Coheight \(0\).} Then \(z\) is a maximal point, hence \(z = \eta_X\), the generic
  point of the irreducible \(X\). But \(\eta_X\) specializes into the dense open
  \(\mathrm{Dom}(f)\), so \(z \in \mathrm{Dom}(f)\), contradicting \(z \in Z(f)\).

  \emph{Coheight \(1\).} Then \(\overline{\{z\}}\) has codimension \(1\), so by
  \cref{lem:smooth_codim_one_dvr} the local ring \(\mathcal O_{X, z}\) is a discrete
  valuation ring with fraction field the function field \(K(X)\). The rational map \(f\)
  determines a morphism \(\operatorname{Spec} K(X) \to Y\) at the generic point, and the
  over-\(\operatorname{Spec}\bar k\) condition \(f \circ Y_{\mathrm{hom}} =
  X_{\mathrm{hom}}\) makes the resulting square commute. Since \(Y\) is proper, the
  valuative criterion of properness (Milne cites Hartshorne~II.4.7) lifts
  \(\operatorname{Spec} K(X) \to Y\) across \(\operatorname{Spec} K(X) \to
  \operatorname{Spec}\mathcal O_{X, z}\) to a morphism \(\operatorname{Spec}
  \mathcal O_{X, z} \to Y\). This spreads to a regular representative of \(f\) on a
  neighbourhood of \(z\), so \(z \in \mathrm{Dom}(f)\), again a contradiction.

  Hence every point of \(Z(f)\) has coheight \(\geq 2\); equivalently \(f\) is
  codim-\(1\)-indeterminacy-free.
\end{proof}

\begin{remark}
  \label{rmk:codim1_role_of_ab}
  \dcref{custom}
  \textbf{Milne 3.1 is valuative-only; Auslander--Buchsbaum is not used here.} The proof
  above uses only normality (via the codim-\(1\) DVR of \cref{lem:smooth_codim_one_dvr})
  and the valuative criterion of properness (Hartshorne~II.4.7), and is
  characteristic-free. There is \emph{no} ``extend across a codim-\(\geq 2\) point'' step:
  the theorem asserts only that \(Z(f)\) has codimension \(\geq 2\), and the actual
  extension of \(f\) to a morphism happens later, once \(Z(f)\) is shown to be empty
  (\cref{thm:codim_one_extension} together with Milne~3.3 gives
  \cref{lem:av_indeterminacyLocus_eq_empty}, \(Z(f) = \emptyset\), and then
  \cref{thm:rational_map_to_av_extends} extends across the empty locus). The
  depth-\(\geq 2\)/local-cohomology (Stacks \texttt{0AVF}) route for a group-variety-free
  target is \emph{not} taken and would in any case be false for a general complete target
  (\emph{cf.}\ \(\mathbb P^2 \dashrightarrow \mathbb P^1\), whose single codim-\(2\)
  indeterminacy point admits no regular extension).

  Where \cref{cor:regular_cohen_macaulay} and the Auslander--Buchsbaum machinery of
  \cref{chap:Albanese_AuslanderBuchsbaum} \emph{do} enter is the Milne~3.3 analysis of a
  group-variety target below (regularity of the diagonal and of local quotients), not
  this codim-\(\geq 2\) statement.
\end{remark}

\section{Milne Lemma 3.3: indeterminacy locus of a map to a group variety is codim 1}
\label{sec:milne_lem33}

The second input to Milne's Theorem~3.2 is structural: when the target \(Y\) is a group
variety, the indeterminacy locus is either empty or of \emph{pure} codimension \(1\).
Combined with \cref{thm:codim_one_extension}'s ``codimension \(\geq 2\)'' conclusion,
this forces the locus to be empty.

The following four lemmas isolate two ingredients of Milne's proof: the
topological existence statement completing the slice argument and Krull's
Hauptidealsatz transported to a scheme-theoretic vanishing statement.  The
first supplies a point where reconstruction applies; the other three show that
a nonzero local function vanishing at a point also vanishes along a
codimension-one generisation.

\begin{lemma}[Slice existence for the fibre of the self-product, and its difference-map corollary]\leanok
  \label{lem:milne33_slice_existence}
  \dcref{custom}
  \lean{AlgebraicGeometry.exists_snd_mem_of_fst_eq_of_mem,
        AlgebraicGeometry.Scheme.RationalMap.exists_mem_domain_precomp_fst_of_differenceRationalMap}
  \uses{lem:difference_map_domain_lower_bound, lem:difference_map_domain_reconstruction}

  \textbf{(Milne Lemma~3.3, Sub-step~2, topological existence half.)} Let \(X\) be
  smooth and geometrically irreducible over \(\bar k\), let \(x \in X\), let
  \(V \subseteq X \times_{\bar k} X\) be an open set meeting the fibre
  \(\mathrm{pr}_1^{-1}\{x\}\) (at some point \(p_0\)), and let \(D \subseteq X\) be a
  dense open subset. Then there is a point \(p\) of the fibre
  \(\mathrm{pr}_1^{-1}\{x\}\) with \(p \in V\) and \(\mathrm{pr}_2(p) \in D\).

  Consequently, if Milne's difference rational map \(\Phi\)
  (\cref{def:rationalMap_difference}) is defined at a point \(p_0\) lying over
  \(x\) --- in particular at the diagonal point \((x, x)\) --- then there is a
  point \(p\) over \(x\), with \(p \in \mathrm{Dom}(\Phi)\), at which the
  precomposition \(f \circ \mathrm{pr}_1\) is defined.
\end{lemma}
\begin{proof}\leanok
  The fibre \(\mathrm{pr}_1^{-1}\{x\}\) is the base change
  \(X_{\kappa(x)}\), hence is irreducible because \(X\) is geometrically
  irreducible. The set \(V\) meets the fibre at \(p_0\) by hypothesis;
  the set \(\mathrm{pr}_2^{-1}(D)\) meets the fibre too. Indeed, for any
  \(u\in D\), the tensor product \(\kappa(x)\otimes_{\bar k}\kappa(u)\)
  is nonzero and therefore has a prime ideal, which determines a point of
  \(X\times_{\bar k}X\) lying over \((x,u)\).
  Preirreducibility of the fibre forces \(V\) and \(\mathrm{pr}_2^{-1}(D)\) to
  meet it \emph{simultaneously}, producing the point \(p\).

  For the corollary, apply the existence statement with \(V := \mathrm{Dom}(\Phi)\)
  (which meets the fibre at \(p_0\)) and \(D := \mathrm{Dom}(f)\), dense since
  \(f\) is a rational map; the resulting point \(p\) satisfies
  \(p \in \mathrm{Dom}(\Phi)\) and \(\mathrm{pr}_2(p) \in \mathrm{Dom}(f)\), so
  \cref{lem:difference_map_domain_reconstruction} places \(p\) in the domain of
  \(f \circ \mathrm{pr}_1\).
\end{proof}

\begin{lemma}[Krull's Hauptidealsatz, transported to a height-one prime below a given prime]\leanok
  \label{lem:krull_hauptidealsatz_transport}
  \dcref{custom}
  \lean{Ideal.exists_height_one_prime_mem_le}

  \textbf{(Milne Lemma~3.3, Sub-step~4, ring-theoretic core.)} Let \(R\) be a
  Noetherian domain, let \(t \in R\) be nonzero, and let \(\mathfrak p\) be a
  prime ideal of \(R\) with \(t \in \mathfrak p\). Then there is a prime ideal
  \(\mathfrak q\) of height exactly \(1\) with \(\mathfrak q \subseteq
  \mathfrak p\) and \(t \in \mathfrak q\).
\end{lemma}
\begin{proof}\leanok
  The principal ideal \((t)\) is contained in \(\mathfrak p\), so it has a
  minimal prime \(\mathfrak q\) lying below \(\mathfrak p\); minimality of
  \(\mathfrak q\) over \((t)\) gives \(t \in \mathfrak q\). Krull's principal
  ideal theorem bounds the height of a minimal prime over a principal ideal
  by \(1\). Since \(R\) is a domain and \(t \neq 0\), the unique height-\(0\)
  prime \((0)\) does not contain \(t\), so \(\mathfrak q \neq (0)\) and hence
  \(\mathrm{height}(\mathfrak q) \geq 1\). The two bounds give
  \(\mathrm{height}(\mathfrak q) = 1\).
\end{proof}

\begin{lemma}[Regularity propagates to a generisation]\leanok
  \label{lem:regular_stalk_propagates_to_generisation}
  \dcref{custom}
  \lean{AlgebraicGeometry.range_algebraMap_stalk_le_of_specializes}

  On an irreducible scheme \(X\), a function-field element regular at a point
  \(P\) is regular at every generisation \(z\) of \(P\) (a point \(z\) with
  \(z \leadsto P\), i.e.\ \(P\) a specialisation of \(z\)): the germ pullback
  \(\mathcal O_{X,P} \to K(X)\) factors through the stalk map
  \(\mathcal O_{X,P} \to \mathcal O_{X,z}\) induced by the specialisation,
  so its image is contained in the image of \(\mathcal O_{X,z} \to K(X)\).
\end{lemma}
\begin{proof}\leanok
  Immediate from functoriality of stalks along the specialisation map: for
  \(s \in \mathcal O_{X,P}\), the germ of \(s\) in the function field equals
  the germ, under \(\mathcal O_{X,z} \to K(X)\), of the image of \(s\) under
  \(\mathcal O_{X,P} \to \mathcal O_{X,z}\).
\end{proof}

\begin{lemma}[Zero loci of regular functions are pure codimension one: the diagonal transport]\leanok
  \label{lem:milne33_diagonal_vanishing_purity}
  \dcref{custom}
  \lean{AlgebraicGeometry.Scheme.exists_specializes_coheight_eq_one_of_mem_maximalIdeal}
  \uses{lem:regular_stalk_propagates_to_generisation}

  \textbf{(Milne Lemma~3.3, Sub-step~4, scheme wrapper.)} Let \(X\) be an
  integral, locally Noetherian scheme with regular local rings at every point
  (e.g.\ \(X\) smooth over a field). Let \(P \in X\) and let \(s\) be an
  element of the maximal ideal of \(\mathcal O_{X,P}\) whose image \(t\) in
  the function field \(K(X)\) is nonzero. Then there is a generisation
  \(z \leadsto P\) of coheight \(1\) and an element \(s'\) of the maximal
  ideal of \(\mathcal O_{X,z}\) whose image in \(K(X)\) is again \(t\).

  Geometrically: the zero locus of a nonzero regular function on \(X\) is of
  pure codimension \(1\), passing through \(P\).
\end{lemma}
\begin{proof}\leanok
  \uses{lem:regular_stalk_propagates_to_generisation}
  Since \(s\) is a non-unit and \(t \neq 0\), the reciprocal \(t^{-1}\) is not
  regular at \(P\): if it were the image of a unit \(u \in \mathcal O_{X,P}\),
  then \(1 - us\) would be a unit mapping to \(0\) in \(K(X)\), which is
  impossible. Purity of the pole locus of the nonzero rational function
  \(t^{-1}\) --- the codimension-\(1\) structure of the zero/pole locus of a
  regular function on a locally Noetherian integral scheme with regular
  stalks, an instance of \cref{lem:krull_hauptidealsatz_transport} applied to
  a local equation of the locus at \(P\) --- produces a coheight-\(1\)
  generisation \(z \leadsto P\) at which \(t^{-1}\) is still not regular. By
  \cref{lem:regular_stalk_propagates_to_generisation}, \(t\) itself remains
  regular at \(z\), witnessed by some \(s' \in \mathcal O_{X,z}\) with image
  \(t\); if \(s'\) were a unit, its inverse would exhibit \(t^{-1}\) as
  regular at \(z\), a contradiction, so \(s'\) lies in the maximal ideal of
  \(\mathcal O_{X,z}\).
\end{proof}

\begin{lemma}[Milne Lemma~3.3: pure-codim-\(1\) structure of the indeterminacy locus into a group variety]\leanok
  \label{lem:milne_codim1_indeterminacy}
  \dcref{custom}
  \lean{AlgebraicGeometry.Scheme.RationalMap.indeterminacy_pure_codim_one_into_grpScheme}
  \uses{def:indeterminacy_locus, def:codim1_cycles, def:order_at_point,
        def:principal_divisor}
  % SOURCE: \cite{milne-av}, Lemma~3.3, \S I.3, p.~17
  % (read from references/abelian-varieties.pdf, PDF page 23).
  % SOURCE QUOTE: "\textsc{Lemma 3.3.} \emph{Let \(\varphi \colon V \dashrightarrow G\)
  % be a rational map from a nonsingular variety to a group variety. Then either
  % \(\varphi\) is defined on all of \(V\) or the points where it is not defined form a
  % closed subset of pure codimension 1 in \(V\) (i.e., a finite union of prime
  % divisors).}"
  \textit{Source: Milne, \emph{Abelian Varieties}, Lemma~3.3, \S I.3, p.~17.}
  Let \(X\) be a nonsingular variety over \(\bar k\), let \(G\) be a group variety over
  \(\bar k\) (a smooth group scheme of finite type), and let \(f \colon X \dashrightarrow
  G\) be a rational map. Then either \(f\) is defined on all of \(X\), or the indeterminacy
  locus \(Z(f) \subseteq X\) is a closed subset of pure codimension \(1\) (i.e.\ a finite
  union of prime divisors).

  Equivalently: it is \emph{never} the case that \(f\) has a non-empty indeterminacy
  locus all of whose irreducible components have codimension \(\geq 2\).
\end{lemma}

% SOURCE QUOTE PROOF: "\textsc{Proof.} Define a rational map
%   \(\Phi \colon V \times V \dashrightarrow G\), $\;(x, y) \mapsto \varphi(x) \cdot
%   \varphi(y)^{-1}$.
% More precisely, if \((U, \varphi_U)\) represents \(\varphi\), then \(\Phi\) is the rational
% map represented by
%   $U \times U \xrightarrow{\varphi_U \times \varphi_U} G \times G
%   \xrightarrow{\mathrm{id} \times \mathrm{inv}} G
%   \xrightarrow{m} G$.
% Clearly \(\Phi\) is defined at a diagonal point \((x, x)\) if \(\varphi\) is defined at
% \(x\), and then \(\Phi(x, x) = e\). Conversely, if \(\Phi\) is defined at \((x, x)\), then
% it is defined on an open neighbourhood of \((x, x)\); in particular, there will be an
% open subset \(U\) of \(V\) such that \(\Phi\) is defined on \(\{x\} \times U\). After
% possible replacing \(U\) by a smaller open subset (not necessarily containing \(x\)),
% \(\varphi\) will be defined on \(U\). For \(u \in U\), the formula
%   \(\varphi(x) = \Phi(x, u) \cdot \varphi(u)\)
% defines \(\varphi\) at \(x\). Thus \(\varphi\) is defined at \(x\) if and only if \(\Phi\) is
% defined at \((x, x)\). The rational map \(\Phi\) defines a map
%   \(\Phi^\ast \colon \mathcal O_{G, e} \to k(V \times V)\).
% Since \(\Phi\) sends \((x, x)\) to \(e\) if it is defined there, it follows that \(\Phi\)
% is defined at \((x, x)\) if and only if
%   \(\operatorname{Im}(\mathcal O_{G, e}) \subset \mathcal O_{V \times V, (x, x)}\).
% Now \(V \times V\) is nonsingular, and so we have a good theory of divisors (AG,
% Chapter 12). For a nonzero rational function \(f\) on \(V \times V\), write
%   \(\operatorname{div}(f) = \operatorname{div}(f)_0 - \operatorname{div}(f)_\infty\),
% with \(\operatorname{div}(f)_0\) and \(\operatorname{div}(f)_\infty\) effective divisors
% --- note that \(\operatorname{div}(f)_\infty = \operatorname{div}(f^{-1})_0\). Then
%   $\mathcal O_{V \times V, (x, x)} = \{f \in k(V \times V) \mid
%   \operatorname{div}(f)_\infty \text{ does not contain } (x, x)\} \cup \{0\}$.
% Suppose \(\phi\) is not defined at \(x\). Then for some $f \in \operatorname{Im}
% (\varphi^\ast)\(, \)(x, x) \in \operatorname{div}(f)_\infty\(, and clearly \)\Phi$ is
% not defined at the points \((y, y) \in \Delta \cap \operatorname{div}(f)_\infty\).
% This is a subset of pure codimension one in \(\Delta\) (AG 9.2), and when we identify
% it with a subset of \(V\), it is a subset of \(V\) of codimension one passing through
% \(x\) on which \(\varphi\) is not defined."
\begin{proof}\leanok
  \uses{def:rationalMap_precomp, lem:rationalMap_precomp_compHom,
        lem:rationalMap_precomp_hom, def:rationalMap_prod, lem:rationalMap_prod_proj,
        def:grpObj_diff, lem:grpObjDiffLeft_comp_hom, def:rationalMap_difference,
        lem:rationalMap_difference_over, lem:compHom_domain_monotone,
        lem:difference_map_domain_lower_bound, lem:rationalMap_precomp_domain_bound,
        lem:rationalMap_prod_domain_bound, lem:grpObj_reconstruction,
        lem:scheme_reconstruction, lem:rationalMap_difference_reconstruct,
        lem:difference_map_domain_reconstruction, lem:milne33_slice_existence,
        lem:milne33_diagonal_vanishing_purity}
  Following Milne's proof of Lemma~3.3. The argument has four sub-steps. Throughout we
  fix a representative \((U, \varphi_U)\) of \(f\).

  The Lean proof takes the component-reduction route to the same conclusion: fix
  \(x \in Z(f)\), choose a Noetherian affine window \(W \ni x\), let \(C\) be the
  irreducible component of \(x\) in \(Z(f) \cap W\), select by Jacobson density a point
  \(x_0 \in C\) closed in \(X\) and avoiding the other components, transport
  sub-step~4 at \(x_0\) to obtain \(z \in Z(f)\) with \(z \rightsquigarrow x_0\) and
  \(\operatorname{coheight} z = 1\), and identify \(z\) with the generic point of \(C\)
  by strictness of coheight along a specialisation. This yields
  \(x \in C = \overline{\{z\}}\).

  \emph{Sub-step 1: the difference rational map.} Define
  \[
    \Phi \colon X \times X \dashrightarrow G,
    \qquad
    \Phi(x, y) \;:=\; \varphi(x) \cdot \varphi(y)^{-1},
  \]
  via the composition
  \[
    U \times U
    \;\xrightarrow{\varphi_U \times \varphi_U}\; G \times G
    \;\xrightarrow{\mathrm{id} \times \mathrm{inv}}\; G \times G
    \;\xrightarrow{m}\; G,
  \]
  where \(\mathrm{inv} \colon G \to G\) is the group-variety inverse and
  \(m \colon G \times G \to G\) is the group law. The map \(\Phi\) is rational by
  composition: precomposing \(f\) with each projection
  \(\mathrm{pr}_i \colon X \times X \to X\) (which is open, being a base change of the
  flat, locally-finitely-presented structure morphism \(X \to \Spec \bar k\)) yields
  rational maps \(f \circ \mathrm{pr}_i \colon X \times X \dashrightarrow G\)
  (\cref{def:rationalMap_precomp}, whose over-\(\bar k\) part is
  \cref{lem:rationalMap_precomp_compHom}, using \cref{lem:rationalMap_precomp_hom} to
  identify \((s_X \circ \mathrm{pr}_i)\) with the structure morphism of \(X \times X\));
  pairing them into \(X \times X \dashrightarrow
  G \times G\) (\cref{def:rationalMap_prod}, whose projection and over-\(\bar k\) properties
  are \cref{lem:rationalMap_prod_proj}) and composing on the right with the difference
  morphism \(\operatorname{diff}_G^{\,\mathrm{left}} = m \circ (\mathrm{id} \times \mathrm{inv})\)
  (\cref{def:grpObj_diff}) gives the named difference map \(\Phi\)
  (\cref{def:rationalMap_difference}); it is again over \(\bar k\)
  (\cref{lem:rationalMap_difference_over}, using \cref{lem:grpObjDiffLeft_comp_hom}). Its
  domain is at least \(U \times U \subseteq X \times X\)
  (\cref{lem:difference_map_domain_lower_bound}), which is dense because \(U\) is
  dense in \(X\).

  \emph{Sub-step 2: \(\Phi\)-defined at \((x, x)\) \(\Leftrightarrow\) \(\varphi\)-defined at
  \(x\).} If \(\varphi\) is defined at \(x\), then \((x,x) \in U \times U \subseteq
  \mathrm{Dom}(\Phi)\) (\cref{lem:difference_map_domain_lower_bound}), so \(\Phi\) is
  defined at \((x, x)\), with value \(\Phi(x, x) = e\) (the identity of \(G\)). Conversely, if \(\Phi\) is
  defined at \((x, x)\), openness of \(\mathrm{Dom}(\Phi)\) in \(X \times X\) produces an
  open set \(W \ni (x, x)\) on which \(\Phi\) is defined; by projection (taking the slice
  \(\{x\} \times W' \subseteq W\) for some open \(W' \subseteq X\), possibly shrinking
  \(W'\) so \(\varphi\) is defined on \(W'\) as well --- \(W'\) need not contain \(x\)), the
  identity
  \[
    \varphi(x) \;=\; \Phi(x, u) \cdot \varphi(u),
    \qquad u \in W',
  \]
  exhibits a regular representative of \(\varphi\) at \(x\). Hence \(x \in
  \mathrm{Dom}(\varphi)\) if and only if \((x, x) \in \mathrm{Dom}(\Phi)\).

  The reconstruction identity \(\varphi(x) = \Phi(x,u) \cdot \varphi(u)\) is
  \cref{lem:rationalMap_difference_reconstruct}, and its domain consequence
  \(\mathrm{Dom}(\Phi) \cap \mathrm{pr}_2^{-1}(\mathrm{Dom}(\varphi)) \subseteq
  \mathrm{Dom}(\varphi \circ \mathrm{pr}_1)\) is
  \cref{lem:difference_map_domain_reconstruction}. A suitable \(u\) exists by openness of
  \(\mathrm{Dom}(\Phi)\) and irreducibility of the fibre \(X_{\kappa(x)}\), which follows
  from geometric irreducibility of \(X\); this is
  \cref{lem:milne33_slice_existence}. Finally, the smooth-descent reflection
  \(\mathrm{Dom}(\varphi \circ \mathrm{pr}_1) \subseteq
  \mathrm{pr}_1^{-1}(\mathrm{Dom}(\varphi))\) transports definedness back to \(X\).

  \emph{Sub-step 3: rephrase via pullback of local rings.} The rational map \(\Phi\)
  pulls \(\mathcal O_{G, e}\) back to a \(\bar k\)-subalgebra of the function field
  \(K(X \times X)\):
  \[
    \Phi^\ast \colon \mathcal O_{G, e} \;\longrightarrow\; K(X \times X).
  \]
  Because \(\Phi\) sends \((x, x)\) to \(e\) wherever it is defined, \(\Phi\) is defined at
  \((x, x)\) if and only if \(\Phi^\ast(\mathcal O_{G, e}) \subseteq \mathcal O_{X
  \times X, (x, x)}\) (i.e.\ every local function on \(G\) near \(e\) has no pole at
  \((x, x)\) after pullback).

  \emph{Sub-step 4: pole locus is pure codimension \(1\).} The variety \(X \times X\) is
  nonsingular (smoothness is stable under products over \(\bar k\)), so the Weil-divisor
  apparatus of \cref{def:codim1_cycles} and the order map of
  \cref{def:order_at_point} apply on \(X \times X\). For any nonzero rational function
  \(f \in K(X \times X)^\ast\) decompose its divisor as
  \[
    \mathrm{div}(f) \;=\; \mathrm{div}(f)_0 - \mathrm{div}(f)_\infty,
  \]
  with \(\mathrm{div}(f)_0\) and \(\mathrm{div}(f)_\infty\) effective and
  \(\mathrm{div}(f)_\infty = \mathrm{div}(f^{-1})_0\) (Milne's notation, matching the
  project's \(\div(f) = \sum \ord_Y(f) \cdot Y\) of
  \cref{def:principal_divisor}: \(\mathrm{div}(f)_0\) collects the prime divisors with
  \(\ord_Y(f) > 0\) and \(\mathrm{div}(f)_\infty\) those with \(\ord_Y(f) < 0\)). The local
  ring at \((x, x)\) can be characterised as
  \[
    \mathcal O_{X \times X, (x, x)}
    \;=\; \bigl\{\, f \in K(X \times X) \;\big|\;
    \mathrm{div}(f)_\infty \text{ does not contain } (x, x) \,\bigr\} \cup \{0\}.
  \]
  Suppose \(f \in K(X \times X)^\ast\) is not in \(\mathcal O_{X \times X, (x, x)}\);
  equivalently, \((x, x) \in \mathrm{supp}(\mathrm{div}(f)_\infty)\). The pole divisor
  \(\mathrm{div}(f)_\infty\) is a finite sum of prime divisors of \(X \times X\) (i.e.\ a
  closed subset of pure codimension \(1\) in \(X \times X\)). Its intersection with the
  diagonal \(\Delta_X \subseteq X \times X\) is a closed subset of pure codimension \(1\)
  in \(\Delta_X\) (this is Milne's reference to AG~9.2, the standard fact that a prime
  divisor of \(X \times X\) either contains \(\Delta_X\) or meets it in pure codimension
  \(1\); on a nonsingular variety this is an instance of Krull's principal-ideal theorem
  applied to the local equation of \(\Delta_X\) inside \(X \times X\) --- the
  scheme-theoretic transport of this purity statement along a generisation of a
  point of \(\Delta_X\) is recorded as \cref{lem:milne33_diagonal_vanishing_purity},
  built from the ring-theoretic core \cref{lem:krull_hauptidealsatz_transport}).

  \emph{Conclusion.} Suppose \(f \colon X \dashrightarrow G\) is not defined on all of
  \(X\). Pick \(x \in Z(f)\). By Sub-steps~2--3, there is some \(f_0 \in
  \Phi^\ast(\mathcal O_{G, e})\) with \((x, x) \in \mathrm{supp}(\mathrm{div}(f_0)_\infty)\).
  By Sub-step~4 the intersection \(\mathrm{supp}(\mathrm{div}(f_0)_\infty) \cap
  \Delta_X\) is pure codimension \(1\) in \(\Delta_X\), passing through \((x, x)\).
  Identifying \(\Delta_X\) with \(X\) via the diagonal isomorphism, the corresponding
  closed subset of \(X\) has pure codimension \(1\) and passes through \(x\). By Sub-step~2,
  every point of this codim-\(1\) subset is a point where \(\Phi\) is undefined on the
  diagonal, hence where \(\varphi\) is undefined. Thus \(Z(\varphi)\) contains a pure
  codim-\(1\) closed subset passing through \(x\); since \(x \in Z(\varphi)\) was arbitrary,
  every irreducible component of \(Z(\varphi)\) has codimension exactly \(1\), i.e.\
  \(Z(\varphi)\) is of pure codimension \(1\).
\end{proof}

\begin{remark}
  \label{rmk:milne_lem33_charfree}
  \dcref{custom}
  \textbf{Characteristic-free.} The above proof uses only: (i) the group law on \(G\)
  and its inverse (formal smoothness of \(G\) as a group object); (ii) the Weil-divisor
  apparatus on the nonsingular variety \(X \times X\) (Hartshorne~II.6, or
  \cref{chap:RiemannRoch_WeilDivisor} in the project's notation); (iii) the
  decomposition of a divisor on a nonsingular variety into effective pieces, and
  Krull's principal-ideal theorem on \(X \times X\). None of these has a characteristic
  restriction, so \cref{lem:milne_codim1_indeterminacy} holds over every
  algebraically closed base field, in arbitrary characteristic. (Milne states the
  lemma in this generality at \S I.3, p.~17.)
\end{remark}

\section{Weil-divisor obstruction to codim-1 extension}
\label{sec:weil_obstruction}

The proof of \cref{thm:codim_one_extension}'s Step~1 (codim-\(1\) ruling-out) goes
through the valuative criterion. In the project's Weil-divisor language
(\cref{chap:RiemannRoch_WeilDivisor}) the same obstruction has a cleaner phrasing:
the extension fails at a prime divisor \(W\) if and only if some affine coordinate of
\(f\) around the image \(\overline{f(\eta_W)}\) has a strictly negative order
\(\ord_W(\cdot) < 0\). This is the valuation-theoretic form of the obstruction
recorded by \cref{def:order_at_point} and \cref{def:principal_divisor}.

\begin{theorem}

  [Weil-divisor criterion for codimension-one extension]
  \label{thm:weil_divisor_obstruction}
  \dcref{custom}
  \lean{AlgebraicGeometry.TODO.weilDivisorObstruction}
  \uses{def:codim_one_indeterminacy, def:order_at_point, lem:smooth_codim_one_dvr,
        lem:mem_domain_partial_map_reshuffle}
  % SOURCE: \cite{hartshorne-ag}, II.6, pp.~130--131 (the valuation \(v_Y\) associated to a
  % prime divisor, and the characterisation of \(\mathcal O_{X, \eta_Y}\) as the
  % subring of \(K(X)\) on which \(v_Y \geq 0\)) (read from
  % references/hartshorne-algebraic-geometry.pdf, PDF pages 147--148).
  % SOURCE QUOTE: "If \(Y\) is a prime divisor on \(X\), let \(\eta \in Y\) be its generic
  % point. Then the local ring \(\mathcal O_{\eta, X}\) of \(X\) at \(\eta\) is a discrete
  % valuation ring with quotient field \(K\), the function field of \(X\). We call the
  % corresponding discrete valuation \(v_Y\) the \emph{valuation of \(Y\)}."
  % SOURCE QUOTE: "Now let \(f \in K^\ast\) be any nonzero rational function on \(X\).
  % Then \(v_Y(f)\) is an integer. If it is positive, we say \(f\) has a \emph{zero} along
  % \(Y\), of that order; if it is negative, we say \(f\) has a \emph{pole} along \(Y\), of
  % order \(-v_Y(f)\)."
  \textit{Source: Hartshorne, II.6, pp.~130--131 (valuation \(v_Y\) and the order map).}
  Let \(X\) be a nonsingular variety over \(\bar k\), let \(Y\) be a variety over \(\bar k\),
  let \(f \colon X \dashrightarrow Y\) be a rational map, and let \(W \subseteq X\) be a
  prime divisor with generic point \(\eta_W \in W\) (so
  \(\mathrm{codim}_X \overline{\{\eta_W\}} = 1\)). Let \((U, \varphi_U)\) be any representative of
  \(f\) with \(U \cap (X \smallsetminus W) \neq \emptyset\). The following are equivalent:
  \begin{enumerate}
    \item \(f\) is defined at \(\eta_W\) (i.e.\ \(\eta_W \in \mathrm{Dom}(f)\), equivalently
      \(W \not\subseteq Z(f)\)).
    \item For some (equivalently, every) affine open neighbourhood \(V \subseteq Y\) of
      the generic image of \(f\) near \(W\), every regular function \(g \in
      \mathcal O_Y(V)\) pulls back to a rational function \(\varphi_U^\ast(g) \in K(X)\)
      satisfying
      \[
        \ord_W\bigl(\varphi_U^\ast(g)\bigr) \;\geq\; 0.
      \]
  \end{enumerate}

  Consequently: \(f\) is codim-\(1\)-indeterminacy-free
  (\cref{def:codim_one_indeterminacy}) if and only if, for every prime divisor
  \(W \subseteq X\) and every affine open \(V \subseteq Y\) as above, the pullback
  \(\varphi_U^\ast(g)\) has \(\ord_W \geq 0\) for all \(g \in \mathcal O_Y(V)\).
\end{theorem}

\begin{proof}

  By \cref{lem:smooth_codim_one_dvr}, the local ring \(\mathcal O_{X, \eta_W}\) is a
  DVR with fraction field \(K(X)\) and uniformiser determined by \(W\); by
  \cref{def:order_at_point} its valuation is exactly \(\ord_W\). The local ring at
  \(\eta_W\) is therefore characterised inside \(K(X)\) as the subring
  \[
    \mathcal O_{X, \eta_W} \;=\; \bigl\{\, h \in K(X) \;\big|\; \ord_W(h) \geq 0
    \,\bigr\} \cup \{0\}.
  \]
  The rational map \(f\) is defined at \(\eta_W\) if and only if, for any affine open
  \(V = \operatorname{Spec} A \subseteq Y\) around the image \(\overline{f(\eta_W)}\),
  the pullback \(\varphi_U^\ast \colon A \to K(X)\) lands in
  \(\mathcal O_{X, \eta_W}\): this is the standard local characterisation of where a
  rational map is regular at a codim-\(1\) point of a smooth variety (it is the
  ``regular = no pole'' criterion). By the displayed identity, the latter is
  equivalent to \(\ord_W(\varphi_U^\ast(g)) \geq 0\) for every \(g \in A\), which is
  condition (2). The equivalence (1)~\(\Leftrightarrow\)~(2) is therefore established.

  For the consequence, \(f\) is codim-\(1\)-indeterminacy-free iff no prime divisor \(W\) is
  contained in \(Z(f)\), which by the equivalence above is iff the \(\ord_W \geq 0\)
  condition holds for every prime divisor \(W\).
\end{proof}

\begin{remark}
  \label{rmk:weil_obstruction_application}
  \dcref{custom}
  \textbf{Role of the valuation form.} In
  \cref{chap:Albanese_Thm32RationalMapExtension}, the abstract codimension statement
  \cref{thm:codim_one_extension} suffices and \cref{thm:weil_divisor_obstruction}
  is not needed. The Weil-divisor form instead isolates an equivalent local
  criterion for regularity at every prime divisor; this criterion is also useful
  for the symmetric-product arguments of \cref{chap:Albanese_AlbaneseUP}.
\end{remark}

\begin{lemma}


  [Definedness at a prime divisor, repackaged as a through-the-point partial map]
  \label{lem:mem_domain_partial_map_reshuffle}
  \dcref{custom}
  Let \(X\) and \(Y\) be varieties over \(\bar k\) and let
  \(f \colon X \dashrightarrow Y\) be a rational map. For a prime divisor
  \(W \subseteq X\) with generic point \(\eta_W = W.\mathrm{point}\), the
  following are equivalent:
  \begin{enumerate}
    \item \(\eta_W \in \mathrm{Dom}(f)\) (\(f\) is defined at the generic point of
      \(W\)).
    \item There exists a partial-map representative \((U, \varphi_U)\) of \(f\) ---
      i.e.\ a regular morphism \(\varphi_U \colon U \to Y\) on a dense open
      \(U \subseteq X\) whose induced rational map equals \(f\) --- whose domain
      \(U\) contains \(\eta_W\).
  \end{enumerate}
\end{lemma}

\begin{proof}

  By definition, \(\eta_W\in\mathrm{Dom}(f)\) precisely when some
  representative of the rational map is defined on an open neighbourhood of
  \(\eta_W\).  Such a representative gives (2), and conversely the domain of
  the representative in (2) witnesses membership in \(\mathrm{Dom}(f)\).
\end{proof}

\begin{remark}
  \label{rmk:mem_domain_partial_map_reshuffle_scope}
  \dcref{custom}
  \textbf{Relationship to \cref{thm:weil_divisor_obstruction}.}
  \cref{lem:mem_domain_partial_map_reshuffle} captures the definitional half of
  the equivalence: a partial-map representative exists through the generic point
  of the prime divisor. The additional mathematical input in
  \cref{thm:weil_divisor_obstruction} is the identification of the local ring with
  the nonnegative part of the associated discrete valuation, applied to the
  pullbacks of regular functions on an affine neighbourhood of the image.
\end{remark}

\section{The codimension-one method}
\label{sec:codim1_lean_encoding}

The results above form a short chain of local and global arguments:
\begin{enumerate}
  \item Smoothness makes every codimension-one local ring a discrete
    valuation ring (\cref{lem:smooth_codim_one_dvr}).
  \item Properness and the valuative criterion then rule out
    codimension-one indeterminacy (\cref{thm:codim_one_extension}).
  \item For a group-variety target, the difference map and its pole
    divisors force every nonempty indeterminacy locus to be pure of
    codimension one (\cref{lem:milne_codim1_indeterminacy}).
  \item For a prime divisor, regularity is equivalently the absence of
    negative order in every pulled-back affine coordinate
    (\cref{thm:weil_divisor_obstruction}).
\end{enumerate}
Thus a rational map from a nonsingular variety to a complete group variety
has empty indeterminacy locus: it has no codimension-one component by (2),
whereas (3) says that every nonempty indeterminacy locus has such a component.

\section{Relations with neighbouring results}
\label{sec:codim1_out_of_scope}

The regular-implies-Cohen--Macaulay theorem
\cref{cor:regular_cohen_macaulay}, obtained from Auslander--Buchsbaum in
\cref{chap:Albanese_AuslanderBuchsbaum}, gives a depth-theoretic description
of the same regular local rings.  The valuative proof of
\cref{thm:codim_one_extension} does not require that stronger description.

Combining \cref{thm:codim_one_extension} and
\cref{lem:milne_codim1_indeterminacy} gives Milne's Theorem~3.2,
\cref{thm:rational_map_to_av_extends}.  Its later rigidity consequences
use the group structure after the extension has been obtained.

The divisor argument needs only the decomposition
\[
  \div(h)=\div(h)_0-\div(h)_\infty
\]
and the codimension-one support of principal divisors.  It does not use the degree of a divisor on a curve or
intersection numbers on a surface.  Cartier divisors could be substituted
for Weil divisors on a smooth variety, where the two theories agree.

The algebraically closed base guarantees that closed points have residue
field \(\bar k\).  Over a general field, the local statements persist after
accounting for residue-field extensions, while descent of the resulting
global morphism is a separate question.
